\section{The pencil in every dimension}
\label{sec:pencil-d}

In projective $d$-space the hinges of a body-hinge framework are
$(d-2)$-flats, and the screw space has dimension $D = \binom{d+1}{2}$.
The hinges at a body form a \emph{pencil} when they lie in a hyperplane,
the panel of the body, and contain a common $(d-3)$-flat of it. At
$d = 3$ this is the pencil of lines through a point in a plane of
\cref{sec:pencil}. At $d = 2$ the $(d-3)$-flat is empty, and a pencil
realization is a hinge-coplanar realization: the body-and-pin frameworks
with collinear pins of Jackson and Jord\'an~\cite{jacksonJordan2008pin}.
A pencil realization satisfies every weaker condition of this kind, the
hinges at each body lying in a hyperplane or containing a common
$j$-flat for some $j \le d - 3$. In particular it is a hinge-coplanar
realization. So a theorem that every multigraph has a pencil realization
at the deficiency rank in dimension $d$ would strengthen Katoh and
Tanigawa's theorem in dimension $d$ (KT
Theorem~5.6~\cite{katohTanigawa2011}), which gives a hinge-coplanar
realization at that rank.

This chapter states pencil realizations in every dimension, and the
objects of \cref{sec:main-component-carrier} that a proof in dimension
$d$ along the lines of \cref{sec:main-component} would use. At $d = 3$ each of
them is the object of \cref{sec:pencil} or \cref{sec:main-component}
that it restates. The third section embeds a lower dimension in a
higher one. At the graphs of small deficiency it carries the attaining
of the general configuration (\cref{thm:pencil-x0-coverage}) from
dimension three to every dimension, and so gives them adjacent-distinct
pencil realizations at the deficiency rank in every dimension
(\cref{sec:pencil-d-stable}). The fourth proves the estimates on
partitions of \cref{sec:main-component-sparse} in every dimension, and
with them a core to contract when no chain is usable
(\cref{sec:pencil-d-deficiency}). The sections after it extend the
induction of \cref{sec:main-component} to every dimension, beginning
with its setting and the rank of a flat configuration
(\cref{sec:pencil-d-induction}).

We describe a pencil realization by its polar. The polarity of
\cref{def:meet-complement-iso} takes the panel with normal $b_v$ to the
point $b_v$, and a hinge in the panels of $u$ and $v$ to a line through
$b_u$ and $b_v$. It takes a $(d-3)$-flat of the panel of $v$ to a plane
through $b_v$, that is, to a three-dimensional subspace $\Pi_v$ of
$K^{d+1}$ containing $b_v$. So the hinges at $v$ form a pencil exactly
when their polar lines pass through $b_v$ and lie in $\Pi_v$: a pencil of
lines through a point in a plane, in every dimension. When the points of
adjacent bodies are independent, the polar of the hinge of an edge $uv$
is the line $b_u \vee b_v$, and the condition says that the points of
each closed neighbourhood lie in a plane
(\cref{lem:pencil-d-point-configuration}).

\subsection{Pencil realizations in every dimension}
\label{sec:pencil-d-realization}

Throughout, $d \ge 2$, and the screw space is $\bigwedge^{d-1} K^{d+1}$.
For a subspace $\Pi$ of $K^{d+1}$ we write $\Pi^{\perp}$ for the vectors
$x$ with $x \cdot y = 0$ for every $y \in \Pi$.

\begin{definition}[Extensor through a subspace]
  \label{def:extensor-through-flat}
  \lean{CombinatorialRigidity.Molecular.ExtensorThroughSubspace}
  \leanok
  \uses{def:screw-space, def:extensor}
  A screw element $C \in \bigwedge^{k} K^{k+2}$ \emph{passes through} a
  subspace $A$ of $K^{k+2}$ when $C$ is the extensor of $k$ vectors of
  $K^{k+2}$ whose span contains $A$.
\end{definition}

\begin{lemma}[Through the span of one vector]
  \label{lem:extensor-through-flat-point}
  \lean{CombinatorialRigidity.Molecular.extensorThroughSubspace_span_singleton_iff}
  \leanok
  \uses{def:extensor-through-flat, def:extensor-through-point}
  A screw element passes through the span of a vector $q$
  (\cref{def:extensor-through-flat}) exactly when it passes through the
  point $q$ (\cref{def:extensor-through-point}).
\end{lemma}
\begin{proof}
  \leanok
  A subspace contains the span of $q$ exactly when it contains $q$.
\end{proof}

\begin{definition}[Pencil realization in every dimension]
  \label{def:pencil-d-realization}
  \lean{CombinatorialRigidity.Molecular.HasPencilPlaneRealization,
        CombinatorialRigidity.Molecular.dotPerp,
        CombinatorialRigidity.Molecular.exists_hasPencilPlaneRealization_one_iff}
  \leanok
  \uses{def:coplanar-panel-realization, def:extensor-through-flat}
  A \emph{pencil realization} of a multigraph $G$ in dimension $d$ is a
  hinge-coplanar panel realization of $G$ in $\bigwedge^{d-1} K^{d+1}$
  (\cref{def:coplanar-panel-realization}), with normal $b_v$ at each
  body $v$, together with a three-dimensional subspace $\Pi_v$ of
  $K^{d+1}$ containing $b_v$ at each body $v$, such that the supporting
  extensor of every link $uv$ passes through $\Pi_u^{\perp}$ and through
  $\Pi_v^{\perp}$ (\cref{def:extensor-through-flat}).
\end{definition}

The subspace $\Pi_v^{\perp}$ has dimension $d - 2$ and lies in the
hyperplane $b_v^{\perp}$: it is a $(d-3)$-flat of the panel of $v$, and
every hinge at $v$ contains it. At $d = 3$ it is spanned by one point,
through which the hinges at $v$ pass. At $d = 2$ it is zero, as $\Pi_v$
is all of $K^3$, so the pencil realizations in dimension two are the
hinge-coplanar panel realizations, with $\Pi_v = K^3$ at each body.

The subspaces $\Pi_v$ are part of the data. At a link whose two normals
are independent, the hinge is determined by them up to scale, and the
condition at the link says that each normal lies in the subspace of the
other end (\cref{lem:pencil-d-point-configuration}). At a link whose two
normals are dependent the panels coincide, and the hinge can be any
$(d-2)$-flat of the common panel. The condition then asks it to contain
the $(d-3)$-flats of both ends. Two bodies joined by two edges reach the
deficiency rank $D$ only at such a link: in two distinct panels both
hinges are multiples of one extensor
(\cref{lem:pencil-d-point-configuration}), and the rank is at most
$D - 1$. If the definition asked only that the normals of each closed
neighbourhood lie in a three-dimensional subspace, then at $d = 3$ two
bodies with one panel, joined by three edges whose hinges are three
lines of that panel with no common point, would form a pencil
realization, though their hinges do not form a pencil.

\begin{remark}[Hinges through a point]
  \label{rem:pencil-d-point-type}
  A pencil realization in the sense of
  \cref{def:pencil-panel-realization} makes sense at every grade: the
  hinges at each body lie in its panel and pass through one point of
  it. For $d \ge 4$ this is weaker than a pencil. In $K^5$ the
  three-dimensional subspaces spanned by $e_1, e_2, e_3$, by
  $e_1, e_2, e_4$ and by $e_1, e_3, e_4$ are planes of projective
  $4$-space that lie in the hyperplane $e_0^{\perp}$ and pass through
  the point $e_1$, but no line lies in all three.
\end{remark}

\begin{lemma}[Pencil realizations in dimension three]
  \label{lem:pencil-d-agrees-three}
  \lean{CombinatorialRigidity.Molecular.hasPencilPlaneRealization_two_iff,
        CombinatorialRigidity.Molecular.exists_hasPencilPlaneRealization_two_iff}
  \leanok
  \uses{def:pencil-d-realization, def:pencil-panel-realization,
        lem:extensor-through-flat-point}
  Let $d = 3$, and let $q_v \in K^4$ be nonzero with
  $\Pi_v = q_v^{\perp}$ at each body $v$. A hinge-coplanar panel
  realization of $G$ with normals $b$ is a pencil realization in
  dimension three with the subspaces $\Pi_v$
  (\cref{def:pencil-d-realization}) exactly when it is a pencil
  realization with normals $b$ and points $q$
  (\cref{def:pencil-panel-realization}). As every three-dimensional
  subspace of $K^4$ is $q^{\perp}$ for a nonzero $q$, it is a pencil
  realization in dimension three for some subspaces exactly when it is a
  pencil realization with normals $b$ for some points.
\end{lemma}
\begin{proof}
  \leanok
  As the dot product of $K^4$ is nondegenerate, $\Pi_v^{\perp}$ is the
  span of $q_v$. So passing through $\Pi_v^{\perp}$ is passing through
  $q_v$ (\cref{lem:extensor-through-flat-point}), and $b_v \in \Pi_v$ is
  the incidence $q_v \cdot b_v = 0$. Every three-dimensional subspace of
  $K^4$ is $q^{\perp}$ for some nonzero $q$.
\end{proof}

\begin{lemma}[Pencil realizations with independent adjacent normals]
  \label{lem:pencil-d-point-configuration}
  \lean{CombinatorialRigidity.Molecular.PanelHingeFramework.toPatchedBodyHinge,
        CombinatorialRigidity.Molecular.PanelHingeFramework.span_rigidityRows_toPatchedBodyHinge,
        CombinatorialRigidity.Molecular.PanelHingeFramework.hasPencilPlaneRealization_toPatchedBodyHinge,
        CombinatorialRigidity.Molecular.HasPencilPlaneRealization.supportExtensor_eq_smul_panelSupportExtensor,
        CombinatorialRigidity.Molecular.HasPencilPlaneRealization.normal_mem_plane_of_mem_closedNbhd}
  \leanok
  \uses{def:pencil-d-realization, def:panel-support-extensor,
        def:panel-hinge-framework, def:rigidity-matrix}
  Let a nonzero $b_v \in K^{d+1}$ and a three-dimensional subspace
  $\Pi_v$ containing $b_v$ be given at each body $v$ of $G$, with $b_u$
  and $b_v$ linearly independent at every edge $uv$.
  \begin{enumerate}
    \item If $b_w \in \Pi_v$ for every body $v$ and every $w$ in the
          closed neighbourhood of $v$, the framework on $G$ whose hinge
          at each edge $uv$ is the panel-support extensor $C(b_u, b_v)$
          (\cref{def:panel-support-extensor}) is, with the normals $b$
          and the subspaces $\Pi$, a pencil realization of $G$ in
          dimension $d$. Its rigidity matrix (\cref{def:rigidity-matrix})
          has the rows of the panel-hinge framework on $G$ with normals
          $b$ (\cref{def:panel-hinge-framework}).
    \item Conversely, in a pencil realization of $G$ in dimension $d$
          with normals $b$ and subspaces $\Pi$, the hinge of every edge
          $uv$ is a nonzero multiple of $C(b_u, b_v)$, and $b_w \in \Pi_v$
          for every $w$ in the closed neighbourhood of every body $v$.
  \end{enumerate}
\end{lemma}
\begin{proof}
  \leanok
  \uses{lem:panel-support-extensor-join, lem:meet-two-hyperplanes,
        lem:decomposable-extensor-span-unique}
  Let $uv$ be an edge. As $b_u$ and $b_v$ are independent, $C(b_u, b_v)$
  is a nonzero extensor of $d - 1$ independent vectors of the
  $(d-1)$-dimensional subspace $b_u^{\perp} \cap b_v^{\perp}$
  (\cref{lem:panel-support-extensor-join}), so of a basis of it. So it
  lies in both panels, and its span contains $\Pi_v^{\perp}$ when $b_u$
  and $b_v$ lie in $\Pi_v$, as then
  $\Pi_v^{\perp} \subseteq b_u^{\perp} \cap b_v^{\perp}$. This gives (1),
  as $u$ and $v$ lie in the closed neighbourhoods of both. Its hinges at
  the edges are those of the panel-hinge framework, so its rows are too.

  For (2), a nonzero extensor determines the span of its vectors, by the
  argument of \cref{lem:decomposable-extensor-span-unique} with $d - 1$
  vectors in place of two. The hinge of $uv$ is nonzero and lies in both
  panels, so it is the extensor of $d - 1$ independent vectors orthogonal
  to $b_u$ and $b_v$, and a nonzero multiple of $C(b_u, b_v)$
  (\cref{lem:meet-two-hyperplanes}). It also passes through
  $\Pi_v^{\perp}$, so $\Pi_v^{\perp}$ lies in the span of vectors
  orthogonal to $b_u$, and
  $b_u \in (\Pi_v^{\perp})^{\perp} = \Pi_v$. This part does not use the
  independence of $b_u$ and $b_v$. As $b_v \in \Pi_v$, it covers the
  closed neighbourhood of $v$, which consists of $v$ and the other ends
  of the links at $v$.
\end{proof}

\begin{fmlnote}
  A body-hinge framework has a hinge at every label of the ambient set
  of edge labels, also at labels that are not edges of $G$. There the
  formal framework of (1) has the fixed nonzero hinge $C(e_0, e_1)$ of
  two standard basis vectors. No rigidity row reads it, and it keeps
  every hinge nonzero, as a hinge-coplanar panel realization requires
  (\cref{def:coplanar-panel-realization}).
\end{fmlnote}

\begin{definition}[Pencil realization at the deficiency rank, in every dimension]
  \label{def:pencil-d-rank-hypothesis}
  \lean{CombinatorialRigidity.Molecular.HasGradedPencilRealization}
  \leanok
  \uses{def:pencil-d-realization, def:D-deficiency, def:rigidity-matrix}
  Let $D = \binom{d+1}{2}$. A multigraph $G$ has a \emph{pencil
  realization at the deficiency rank} in dimension $d$ when it has a
  pencil realization in dimension $d$ (\cref{def:pencil-d-realization})
  whose rigidity rows span a space of dimension
  $D(|V(G)| - 1) - \operatorname{def}(\tilde G)$, the deficiency taken
  at $D$ (\cref{def:D-deficiency}).
\end{definition}

\begin{definition}[Adjacent-distinct pencil realization, in every dimension]
  \label{def:pencil-d-distinct-motive}
  \lean{CombinatorialRigidity.Molecular.HasGradedDistinctPencilRealization}
  \leanok
  \uses{def:pencil-d-rank-hypothesis}
  A multigraph $G$ has an \emph{adjacent-distinct} pencil realization
  in dimension $d$ when it has a pencil realization at the deficiency
  rank in dimension $d$ (\cref{def:pencil-d-rank-hypothesis}) whose
  normals $b_u$ and $b_v$ are linearly independent at every link $uv$.
\end{definition}

In the polar form the points of adjacent bodies are distinct, as in
\cref{def:pencil-distinct-motive}. By
\cref{lem:pencil-d-point-configuration}, an adjacent-distinct pencil
realization in dimension $d$ is determined by its normals up to the
scale of each hinge, which does not change the rank.

\begin{lemma}[At the deficiency rank in dimension three]
  \label{lem:pencil-d-rank-agrees-three}
  \lean{CombinatorialRigidity.Molecular.hasGradedPencilRealization_two_iff,
        CombinatorialRigidity.Molecular.hasGradedDistinctPencilRealization_two_iff}
  \leanok
  \uses{def:pencil-d-rank-hypothesis, def:pencil-d-distinct-motive,
        def:pencil-rank-hypothesis, def:pencil-distinct-motive}
  A multigraph has a pencil realization at the deficiency rank in
  dimension three (\cref{def:pencil-d-rank-hypothesis}) exactly when it
  has a pencil realization at the deficiency rank
  (\cref{def:pencil-rank-hypothesis}). It has an adjacent-distinct one
  in dimension three (\cref{def:pencil-d-distinct-motive}) exactly when
  it has an adjacent-distinct pencil realization
  (\cref{def:pencil-distinct-motive}).
\end{lemma}
\begin{proof}
  \leanok
  \uses{lem:pencil-d-agrees-three, lem:pencil-self-dual, lem:screw-map-rows}
  The first equivalence is \cref{lem:pencil-d-agrees-three}, applied to
  one framework, so at one rank. For the second, the adjacent-distinct
  realizations in dimension three have independent normals at the ends
  of every link, and those of \cref{def:pencil-distinct-motive} have
  independent points there. The polarity exchanges the normals and the
  points of a pencil realization (\cref{lem:pencil-self-dual}), and it
  keeps the rank (\cref{lem:screw-map-rows}).
\end{proof}

\subsection{Configurations in every dimension}
\label{sec:pencil-d-carrier}

A configuration in dimension $d$ lifts a planar picture by $d - 2$
heights, each of them in the lifting space $L(q)$ of
\cref{def:pencil-lifting-space}. The rank of a configuration is read at
the panel-hinge framework whose normals are its points, as at $d = 3$
(\cref{def:pencil-x0-attains}).

\begin{definition}[Configurations in every dimension]
  \label{def:pencil-d-configuration}
  \lean{CombinatorialRigidity.Molecular.pencilConfigPointGrade,
        CombinatorialRigidity.Molecular.pencilConfigPointGrade_eq_pencilPicturePoint,
        CombinatorialRigidity.Molecular.pencilConfigPointGrade_eq_pencilConfigPoint,
        CombinatorialRigidity.Molecular.PanelHingeFramework.finrank_span_rigidityRows_ofNormals_pencilConfigPointGrade_one}
  \leanok
  \uses{def:pencil-lifting-space, def:panel-hinge-framework,
        def:rigidity-matrix, def:pencil-admissible-picture,
        def:pencil-configuration}
  A planar picture $q$ of $G$ and heights
  $z^1, \dots, z^{d-2} \colon V(G) \to K$ give the
  \emph{configuration} $b = (q, z)$ of the points
  $b_v = (x_v, y_v, z^1_v, \dots, z^{d-2}_v, 1) \in K^{d+1}$. Its
  \emph{rank} is the rank of the rigidity matrix
  (\cref{def:rigidity-matrix}) of the panel-hinge framework on $G$ in
  $\bigwedge^{d-1} K^{d+1}$ whose panel normal at each body $v$ is $b_v$
  (\cref{def:panel-hinge-framework}).
\end{definition}

At $d = 3$ these are the configurations of \cref{def:pencil-configuration}
and their rank. At $d = 2$ there is no height, and $b_v$ is the
homogeneous picture point $(x_v, y_v, 1)$ of
\cref{def:pencil-admissible-picture}. If every $z^i$ lies in $L(q)$,
each height is affine on each closed neighbourhood, so the points of
each closed neighbourhood lie in a plane.

\begin{lemma}[A configuration over an admissible picture is a pencil realization]
  \label{lem:pencil-d-config-realization}
  \lean{Graph.IsAdmissiblePicture.hasPencilPlaneRealization_ofNormals_pencilConfigPointGrade,
        Graph.IsAdmissiblePicture.hasGradedDistinctPencilRealization,
        CombinatorialRigidity.Molecular.linearIndependent_pencilConfigPointGrade_pair,
        CombinatorialRigidity.Molecular.linearIndependent_pencilConfigPointGrade_of_linearIndependent_pencilPicturePoint,
        CombinatorialRigidity.Molecular.finrank_span_pencilConfigPointGrade_closedNbhd_le}
  \leanok
  \uses{def:pencil-d-configuration, def:pencil-admissible-picture,
        def:pencil-d-distinct-motive, lem:pencil-d-point-configuration}
  Let $q$ be an admissible picture of $G$, let $z^1, \dots, z^{d-2}$ lie
  in $L(q)$, and let $\Pi_v$ be the span of the points of the closed
  neighbourhood of $v$ in the configuration $b = (q, z)$. Then the
  panel-hinge framework of \cref{def:pencil-d-configuration}, with the
  subspaces $\Pi_v$, is a pencil realization of $G$ in dimension $d$
  with independent normals at the ends of every link. If $(q, z)$ has
  rank $D(|V(G)| - 1) - \operatorname{def}(\tilde G)$, the deficiency
  taken at $D = \binom{d+1}{2}$, then $G$ has an adjacent-distinct pencil
  realization in dimension $d$.
\end{lemma}
\begin{proof}
  \leanok
  On the closed neighbourhood of $v$ each $z^i$ agrees with an affine
  function $a^i x + c^i y + e^i$ of the picture. So the points $b_w$ there
  are the images of the homogeneous picture points $(x_w, y_w, 1)$ under
  one linear map $K^3 \to K^{d+1}$, injective as the first, second and
  last coordinates are kept. The picture points of the closed
  neighbourhood are not collinear, so $\Pi_v$ has dimension three.
  Adjacent picture points differ, so adjacent points $b_u$ and $b_v$ are
  independent, as in \cref{lem:pencil-config-point-join-rank}. Then
  \cref{lem:pencil-d-point-configuration} applies.
\end{proof}

\begin{definition}[The general configuration attains, in every dimension]
  \label{def:pencil-d-x0-attains}
  \lean{Graph.X0AttainsGrade, Graph.liftingSpaceGrade, Graph.mem_liftingSpaceGrade}
  \leanok
  \uses{def:pencil-d-configuration, def:pencil-admissible-picture,
        def:pencil-lifting-space, def:D-deficiency}
  Let $D = \binom{d+1}{2}$. The \emph{general configuration of $G$
  attains} in dimension $d$ if there is a nonzero polynomial $P$ in the
  picture coordinates such that every picture $q$ with $P(q) \neq 0$ is
  admissible and the heights $z \in L(q)^{d-2}$ at which $(q, z)$ has
  rank $D(|V(G)| - 1) - \operatorname{def}(\tilde G)$, the deficiency
  taken at $D$, contain a nonempty Zariski-open subset of $L(q)^{d-2}$.
\end{definition}

\begin{lemma}[Attaining in dimension three]
  \label{lem:pencil-d-x0-attains-three}
  \lean{Graph.x0AttainsGrade_one_iff}
  \leanok
  \uses{def:pencil-d-x0-attains, def:pencil-x0-attains}
  The general configuration of $G$ attains in dimension three
  (\cref{def:pencil-d-x0-attains}) exactly when it attains in the sense
  of \cref{def:pencil-x0-attains}.
\end{lemma}
\begin{proof}
  \leanok
  At $d = 3$ there is one height, and the two definitions read the same
  rank of the same framework.
\end{proof}

\begin{lemma}[A collineation of the normals preserves the rank, in every dimension]
  \label{lem:pencil-d-rank-collineation}
  \lean{CombinatorialRigidity.Molecular.PanelHingeFramework.finrank_span_rigidityRows_ofNormals_linearEquiv_grade,
        CombinatorialRigidity.Molecular.panelSupportExtensor_linearEquiv_eq_det_smul}
  \leanok
  \uses{def:panel-hinge-framework, def:rigidity-matrix}
  Let $g$ be an invertible linear map of $K^{d+1}$. If panel normals $b$
  and $b'$ satisfy $b'_w = g(b_w)$ at every body $w$ of $G$, the
  rigidity matrices of the panel-hinge frameworks on $G$ in
  $\bigwedge^{d-1} K^{d+1}$ with normals $b$ and with normals $b'$ have
  the same rank.
\end{lemma}
\begin{proof}
  \leanok
  \uses{def:panel-support-extensor, lem:complement-iso-toDual,
        def:induced-screw-map, lem:screw-map-rows}
  The hinge $C(b_u, b_v)$ is the image of the join $b_u \vee b_v$ under
  the complement isomorphism (\cref{def:panel-support-extensor}). By its
  defining property (\cref{lem:complement-iso-toDual}), that isomorphism
  takes the map induced by $g$ on $\bigwedge^2 K^{d+1}$ to $\det g$ times
  the screw map $\Lambda_h$ (\cref{def:induced-screw-map}) of the inverse
  transpose $h$ of $g$, as in the proof of
  \cref{lem:meet-two-hyperplanes}. So
  $C(g b_u, g b_v)$ is the
  image of $C(b_u, b_v)$ under one invertible linear map of the screw
  space, the same at every edge, and the two ranks agree
  (\cref{lem:screw-map-rows}).
\end{proof}

\begin{lemma}[Linear changes and affine shifts of the heights preserve the rank]
  \label{lem:pencil-d-rank-scale-shift}
  \lean{Graph.finrank_span_rigidityRows_ofNormals_heights_linearEquiv_add_affineLifts}
  \leanok
  \uses{def:pencil-d-configuration, def:pencil-lifting-space}
  Let $q$ be a picture of $G$ and $z = (z^1, \dots, z^{d-2})$ heights.
  Let $T$ be an invertible $(d-2) \times (d-2)$ matrix over $K$ and
  $a^1, \dots, a^{d-2} \in \mathrm{Aff}(q)$. Then the configurations
  $(q, z)$ and $(q, Tz + a)$ have the same rank, where
  $(Tz + a)^i = \sum_j T_{ij} z^j + a^i$.
\end{lemma}
\begin{proof}
  \leanok
  \uses{lem:pencil-d-rank-collineation}
  Write $a^i_w = \alpha_i x_w + \beta_i y_w + \gamma_i$, and split
  $K^{d+1} = K^2 \times K^{d-2} \times K$. The linear map
  $(\xi, \eta, \zeta, \omega) \mapsto
  (\xi, \eta, T\zeta + \xi\alpha + \eta\beta + \omega\gamma, \omega)$ is
  invertible, since $T$ is. It maps the point of $(q, z)$ at each body to
  the point of $(q, Tz + a)$ there. So
  \cref{lem:pencil-d-rank-collineation} applies.
\end{proof}

\begin{lemma}[One attaining configuration forces the general one, in every dimension]
  \label{lem:pencil-d-x0-one-witness}
  \lean{Graph.x0AttainsGrade_of_exists}
  \leanok
  \uses{def:pencil-main-picture, def:pencil-d-x0-attains,
        lem:rank-polynomial-of-le-finrank,
        lem:relative-deficiency-rank-bound}
  Let $G$ have at least one body, let $q_0$ be a main picture of $G$
  (\cref{def:pencil-main-picture}) and $z_0 \in L(q_0)^{d-2}$, and
  suppose the configuration $(q_0, z_0)$ has rank at least
  $D(|V(G)| - 1) - \operatorname{def}(\tilde G)$, the deficiency taken
  at $D = \binom{d+1}{2}$. Then the general configuration of $G$ attains
  in dimension $d$.
\end{lemma}
\begin{proof}
  \leanok
  \uses{def:pencil-lifting-system, lem:pencil-polynomial-section}
  The proof of \cref{lem:pencil-x0-one-witness} applies, with the
  polynomial section of the lifting system
  (\cref{lem:pencil-polynomial-section}) taken once for each height: it
  gives heights $Z^i(q) \in L(q)$, polynomial in the picture $q$, with
  $Z^i(q_0) = z^i_0$. Over an admissible picture adjacent points are
  independent, so every hinge is nonzero, and the rank polynomial of
  \cref{lem:rank-polynomial-of-le-finrank} is taken in the coordinates of
  the points of $K^{d+1}$. The rank never exceeds the target
  (\cref{lem:relative-deficiency-rank-bound}).
\end{proof}

\begin{lemma}[The rank of the core along the curve, in every dimension]
  \label{lem:pencil-d-contract-magnified-rank}
  \lean{Graph.finrank_span_rigidityRows_induce_contractHeightGrade_eq}
  \leanok
  \uses{def:pencil-contract-lifting-system, def:pencil-d-configuration}
  Let $r \in W$ and $H = G[W]$, let $q$ be a picture and $t \in K$
  nonzero, and for each height $i$ let $(z^i, h^i)$ be a vector of
  unknowns of the rescaled lifting system
  (\cref{def:pencil-contract-lifting-system}). Then the rank of $H$ at
  the configuration $(q(t), z(t))$ in dimension $d$, where
  $z(t) = (z^1(t), \dots, z^{d-2}(t))$, equals its rank at $(q, z)$.
\end{lemma}
\begin{proof}
  \leanok
  \uses{lem:pencil-d-rank-collineation}
  At a body $c$ of $W$ the point of $(q(t), z(t))$ is
  $((1 - t)\,x_r + t\,x_c,\ (1 - t)\,y_r + t\,y_c,\
  t\,z^1_c, \dots, t\,z^{d-2}_c,\ 1)$, the image of the point of
  $(q, z)$ under the linear map
  $(\xi, \eta, \zeta, \omega) \mapsto
  (t\,\xi + (1 - t)\,x_r\,\omega,\ t\,\eta + (1 - t)\,y_r\,\omega,\
  t\,\zeta,\ \omega)$ of $K^{d+1}$, which is invertible as $t \neq 0$. So
  \cref{lem:pencil-d-rank-collineation} applies.
\end{proof}

\begin{lemma}[Independence of joins is an open condition, in every dimension]
  \label{lem:pencil-d-join-independence-open}
  \lean{CombinatorialRigidity.Molecular.exists_mvPolynomial_linearIndependent_join_picture_grade,
        CombinatorialRigidity.Molecular.exists_mvPolynomial_linearIndependent_join_heights_grade}
  \leanok
  \uses{def:pencil-d-configuration, def:extensor}
  Let $(u_i, w_i)_{i \in I}$ be a finite family of pairs of bodies, and
  write $b$ for the points of a configuration $(q, z)$ in dimension $d$.
  If the joins $b_{u_i} \vee b_{w_i} \in \bigwedge^2 K^{d+1}$ are
  linearly independent at $(q_0, z)$, there is a polynomial $P$ in the
  picture, nonzero at $q_0$, such that they are linearly independent at
  $(q, z)$ whenever $P(q) \neq 0$. Likewise, if they are independent at
  $(q, z_0)$, there is a polynomial $R$ in the heights, nonzero at $z_0$,
  such that they are independent at $(q, z)$ whenever $R(z) \neq 0$.
\end{lemma}
\begin{proof}
  \leanok
  The coordinates of $b_u \vee b_w$ are the $2 \times 2$ minors of the
  two points, polynomials in the picture at fixed heights and in the
  heights at a fixed picture. A finite family of vectors is linearly
  independent exactly when some maximal minor of its coordinate matrix
  is nonzero; take for $P$, or $R$, one such minor.
\end{proof}

\subsection{From a lower dimension}
\label{sec:pencil-d-stable}

A pencil realization in dimension $d'$ becomes one in dimension
$d \ge d'$ under the inclusion of $K^{d'+1}$ in $K^{d+1}$ that inserts
zero coordinates before the last one; at a configuration it appends zero
heights. On a connected multigraph the rank grows by
$(D - D')(|V(G)| - 1)$, where $D' = \binom{d'+1}{2}$ and
$D = \binom{d+1}{2}$. Write $\operatorname{def}_d(G)$ for the deficiency
at $D$ (\cref{def:D-deficiency}), as before
\cref{lem:deficiency-antitone}. Then a realization at the deficiency
rank in dimension $d'$ gives one in dimension $d$ when
$\operatorname{def}_{d'}(G) = \operatorname{def}_d(G)$, and on a
connected multigraph the deficiency stops changing once $D$ is large
compared with $\operatorname{def}_2(G)$. So the general configuration,
which attains in dimension three at every graph satisfying the standing
hypotheses (\cref{thm:pencil-x0-coverage}), attains in every dimension
at those with $\operatorname{def}_2(G) \le 3$
(\cref{thm:pencil-d-deficiency-two-le-three}).

\begin{lemma}[The deficiency stops changing]
  \label{lem:deficiency-stable}
  \lean{Graph.Connected.deficiency_eq_of_deficiency_two_add_three_le,
        Graph.Connected.deficiency_eq_of_two_add_deficiency_two_le}
  \leanok
  \uses{def:D-deficiency}
  Let $G$ be a connected multigraph, and let $d' \le d$ with
  $\operatorname{def}_2(G) + 3 \le \binom{d'+1}{2}$, which holds in
  particular when $\operatorname{def}_2(G) + 2 \le d'$. Then
  $\operatorname{def}_{d'}(G) = \operatorname{def}_d(G)$.
\end{lemma}
\begin{proof}
  \leanok
  \uses{lem:deficiency-antitone}
  For a partition $P$ of $V(G)$ let $c(P) = d_G(P) - (|P| - 1)$, which
  is nonnegative as $G$ is connected (the proof of
  \cref{lem:deficiency-antitone}). The value of $P$ at $D$ is
  $D(|P| - 1) - (D - 1)\,d_G(P) = |P| - 1 - (D - 1)\,c(P)$. At $D = 3$
  it is at most $t = \operatorname{def}_2(G)$, so at $D \ge 3$ it is at
  most $t - (D - 3)\,c(P)$. Let $D' = \binom{d'+1}{2}$, so that
  $D' \ge t + 3$, and let $D \ge D'$. If $c(P) \ge 1$, the value of $P$
  at $D$ is at most $t - (D' - 3) \le 0$. If $c(P) = 0$, it is
  $|P| - 1$ at every $D$, and the partition $\{V(G)\}$ has value $0$.
  So at every $D \ge D'$, in particular at $D'$ and at
  $\binom{d+1}{2}$, the deficiency is the largest $|P| - 1$ over the
  partitions with $c(P) = 0$, which does not depend on $D$. If
  $t + 2 \le d'$, then $D' \ge \binom{t+3}{2} \ge t + 3$, as $t \ge 0$.
\end{proof}

\begin{definition}[Embedding a lower dimension]
  \label{def:pencil-d-embedding}
  \lean{CombinatorialRigidity.Molecular.gradeEmbed,
        CombinatorialRigidity.Molecular.screwEmbed,
        CombinatorialRigidity.Molecular.gradeEmbed_injective,
        CombinatorialRigidity.Molecular.screwEmbed_injective}
  \leanok
  \uses{def:screw-space, def:meet-complement-iso, def:join}
  Let $2 \le d' \le d$. The \emph{embedding}
  $\iota \colon K^{d'+1} \to K^{d+1}$ inserts $d - d'$ zero coordinates
  before the last one:
  \[
    \iota(x_1, \dots, x_{d'}, x_{d'+1})
      = (x_1, \dots, x_{d'}, 0, \dots, 0, x_{d'+1}).
  \]
  It induces the linear map $\bigwedge^2 \iota$ of
  $\bigwedge^2 K^{d'+1}$ into $\bigwedge^2 K^{d+1}$ that takes each join
  $u \vee w$ to $\iota u \vee \iota w$. The \emph{screw embedding}
  $\varphi \colon \bigwedge^{d'-1} K^{d'+1} \to \bigwedge^{d-1} K^{d+1}$
  is the inverse of the complement isomorphism of $K^{d'+1}$
  (\cref{def:meet-complement-iso}), followed by $\bigwedge^2 \iota$ and
  the complement isomorphism of $K^{d+1}$. Both $\iota$ and $\varphi$
  are linear and injective.
\end{definition}
\begin{proof}
  \leanok
  The map $\iota$ keeps every coordinate of its argument, and the map
  $\pi$ that deletes the inserted coordinates is a linear left inverse
  of it. Then $\bigwedge^2 \pi$ is a left inverse of $\bigwedge^2 \iota$,
  so $\bigwedge^2 \iota$ is injective, and $\varphi$ is its composite
  with two isomorphisms.
\end{proof}

\begin{lemma}[The screw embedding of a hinge of two panels]
  \label{lem:pencil-d-embedding-join}
  \lean{CombinatorialRigidity.Molecular.screwEmbed_panelSupportExtensor,
        CombinatorialRigidity.Molecular.exteriorPower_map_normalsJoin}
  \leanok
  \uses{def:pencil-d-embedding, def:panel-support-extensor}
  For $u, w \in K^{d'+1}$, the screw embedding takes the panel-support
  extensor $C(u, w)$ (\cref{def:panel-support-extensor}) to
  $C(\iota u, \iota w)$.
\end{lemma}
\begin{proof}
  \leanok
  The extensor $C(u, w)$ is the image of the join $u \vee w$ under the
  complement isomorphism (\cref{def:panel-support-extensor}). So
  $\varphi(C(u, w))$ is the image of
  $(\bigwedge^2 \iota)(u \vee w) = \iota u \vee \iota w$ under the
  complement isomorphism of $K^{d+1}$, which is
  $C(\iota u, \iota w)$.
\end{proof}

\begin{lemma}[An injective map of screw spaces adds to the rank]
  \label{lem:screw-embedding-rank}
  \lean{CombinatorialRigidity.Molecular.BodyHingeFramework.map,
        CombinatorialRigidity.Molecular.BodyHingeFramework.finrank_span_rigidityRows_map_injective,
        CombinatorialRigidity.Molecular.BodyHingeFramework.finrank_infinitesimalMotions_map_injective}
  \leanok
  \uses{def:hinge-constraint, def:rigidity-matrix}
  Let $d' \le d$, let $\varphi$ be an injective linear map of
  $\bigwedge^{d'-1} K^{d'+1}$ into $\bigwedge^{d-1} K^{d+1}$, and let
  $(G, p)$ be a body-hinge framework in $K^{d'}$ on a connected
  multigraph $G$. Let $(G, \varphi p)$ be the body-hinge framework in
  $K^d$ on $G$ whose supporting extensor at each edge $e$ is
  $\varphi\, C(p(e))$. Then
  \[
    \operatorname{rank} R(G, \varphi p)
      = \operatorname{rank} R(G, p) + (D - D')(|V(G)| - 1),
  \]
  where $D' = \binom{d'+1}{2}$ and $D = \binom{d+1}{2}$.
\end{lemma}
\begin{proof}
  \leanok
  Let $B \supseteq V(G)$ be the set of bodies and $m = |B \setminus V(G)|$.
  Write $\varphi S$ for the screw assignment $u \mapsto \varphi(S(u))$,
  and let $Z_0$ and $Z_0'$ be the screw assignments in dimensions $d$ and
  $d'$ that take the same value at the two ends of every edge. As $G$ is
  connected, these are the assignments constant on $V(G)$, so
  $\dim Z_0 = D(1 + m)$ and $\dim Z_0' = D'(1 + m)$. As $\varphi$ is
  injective, $\varphi(x)$ lies in the span of $\varphi\, C(p(e))$
  exactly when $x$ lies in the span of $C(p(e))$. So $\varphi Z(G, p)$
  and $Z_0$ lie in $Z(G, \varphi p)$. Conversely, let
  $S \in Z(G, \varphi p)$ and $v_0 \in V(G)$. Each hinge constraint of
  $(G, \varphi p)$ asks the difference of the values at the ends of an
  edge to lie in the image of $\varphi$, so, as $G$ is connected,
  $S(u) - S(v_0)$ lies in the image of $\varphi$ for every $u \in V(G)$.
  Let $S_0 \in Z_0$ take the value $S(v_0)$ on $V(G)$ and agree with $S$
  off it. Then $S - S_0 = \varphi S'$ for a screw assignment $S'$, which
  lies in $Z(G, p)$, as $\varphi(S'(u) - S'(v)) = S(u) - S(v)$ at each
  edge with ends $u$ and $v$. So $Z(G, \varphi p) = \varphi Z(G, p) + Z_0$.
  An assignment $\varphi S'$ lies in $Z_0$ exactly when $S' \in Z_0'$,
  and $Z_0' \subseteq Z(G, p)$, so the intersection of the two is
  $\varphi Z_0'$. Hence
  $\dim Z(G, \varphi p) + D'(1 + m) = \dim Z(G, p) + D(1 + m)$. The ranks
  are $D|B| - \dim Z(G, \varphi p)$ and $D'|B| - \dim Z(G, p)$
  (\cref{def:rigidity-matrix}), and $|B| = |V(G)| + m$.
\end{proof}

\begin{lemma}[A pencil realization in a lower dimension]
  \label{lem:pencil-d-embed-realization}
  \lean{CombinatorialRigidity.Molecular.HasPencilPlaneRealization.map_gradeEmbed,
        CombinatorialRigidity.Molecular.hasGradedPencilRealization_of_le,
        CombinatorialRigidity.Molecular.hasGradedDistinctPencilRealization_of_le}
  \leanok
  \uses{def:pencil-d-embedding, def:pencil-d-realization,
        def:pencil-d-rank-hypothesis, def:pencil-d-distinct-motive,
        def:D-deficiency}
  Let $2 \le d' \le d$, with $D'$ and $D$ as in
  \cref{lem:screw-embedding-rank}. Let a pencil realization of $G$ in
  dimension $d'$ (\cref{def:pencil-d-realization}) have hinges $C_e$,
  normals $b_v$ and subspaces $\Pi_v$. Then the framework on $G$ with the
  hinges $\varphi(C_e)$, the normals $\iota b_v$ and the subspaces
  $\iota \Pi_v$ (\cref{def:pencil-d-embedding}) is a pencil realization
  of $G$ in dimension $d$, and its normals are independent at every link
  where the given ones are. If $G$ is connected and the given rigidity
  rows span a space of dimension $r$, the new ones span a space of
  dimension $r + (D - D')(|V(G)| - 1)$. So, for $G$ connected and every
  $d''$, a pencil realization of $G$ in dimension $d'$ whose rigidity
  rows span a space of dimension
  $D'(|V(G)| - 1) - \operatorname{def}_{d''}(G)$ gives one in dimension
  $d$ whose rows span a space of dimension
  $D(|V(G)| - 1) - \operatorname{def}_{d''}(G)$, with independent
  normals at every link if the given one has them.
\end{lemma}
\begin{proof}
  \leanok
  \uses{lem:screw-embedding-rank, lem:pencil-d-embedding-join,
        lem:pencil-d-point-configuration, lem:panel-support-extensor-join}
  Let $e$ be a link $uv$. Its hinge $C_e$ is a nonzero extensor of
  $d' - 1$ vectors spanning a subspace $A$ of $K^{d'+1}$, so $A^{\perp}$
  has dimension two. That $C_e$ lies in the panel of $v$ says
  $b_v \in A^{\perp}$, and that it passes through $\Pi_v^{\perp}$ says
  $\Pi_v^{\perp} \subseteq A$, that is, $A^{\perp} \subseteq \Pi_v$;
  likewise at $u$. Let $w_1, w_2$ be a basis of $A^{\perp}$. Then $C_e$
  lies in the panels of $w_1$ and $w_2$, so it is a nonzero multiple of
  $C(w_1, w_2)$, by the argument of part~(2) of
  \cref{lem:pencil-d-point-configuration}. So $\varphi(C_e)$ is the same
  multiple of $C(\iota w_1, \iota w_2)$
  (\cref{lem:pencil-d-embedding-join}), a nonzero extensor of vectors
  spanning the orthogonal of $\iota(A^{\perp})$
  (\cref{lem:panel-support-extensor-join}). As
  $\iota b_v \in \iota(A^{\perp}) \subseteq \iota \Pi_v$, the hinge
  $\varphi(C_e)$ lies in the panel of $\iota b_v$ and passes through
  $(\iota \Pi_v)^{\perp}$. As $\iota$ is injective, each $\iota \Pi_v$
  has dimension three and contains $\iota b_v \neq 0$, and independent
  normals stay independent; as $\varphi$ is injective, every hinge
  $\varphi(C_e)$ is nonzero. The rows are
  \cref{lem:screw-embedding-rank}. For the last claim,
  $D'(|V(G)| - 1) - \operatorname{def}_{d''}(G) + (D - D')(|V(G)| - 1)
  = D(|V(G)| - 1) - \operatorname{def}_{d''}(G)$.
\end{proof}

With $d'' = d'$: if $\operatorname{def}_{d'}(G) = \operatorname{def}_d(G)$
and the connected multigraph $G$ has a pencil realization at the
deficiency rank in dimension $d'$ (\cref{def:pencil-d-rank-hypothesis}),
or an adjacent-distinct one (\cref{def:pencil-d-distinct-motive}), then
it has one in dimension $d$.

\begin{lemma}[Zero heights add to the rank]
  \label{lem:pencil-d-embed-config-rank}
  \lean{CombinatorialRigidity.Molecular.padHeights,
        CombinatorialRigidity.Molecular.pencilConfigPointGrade_padHeights,
        Graph.finrank_span_rigidityRows_ofNormals_pencilConfigPointGrade_padHeights,
        Graph.padHeights_mem_liftingSpaceGrade}
  \leanok
  \uses{def:pencil-d-configuration, def:pencil-d-embedding,
        def:pencil-lifting-space}
  Let $G$ be connected and $2 \le d' \le d$, with $D'$ and $D$ as in
  \cref{lem:screw-embedding-rank}. Let $q$ be a picture of $G$ and
  $z' = (z^1, \dots, z^{d'-2})$ heights, and let
  $z = (z^1, \dots, z^{d'-2}, 0, \dots, 0)$ append $d - d'$ zero
  heights. Then the rank of $(q, z)$ in dimension $d$ is the rank of
  $(q, z')$ in dimension $d'$ plus $(D - D')(|V(G)| - 1)$. If every
  $z^i$ lies in $L(q)$, so does every height of $z$.
\end{lemma}
\begin{proof}
  \leanok
  \uses{lem:pencil-d-embedding-join, lem:screw-embedding-rank}
  The point of $(q, z)$ at each body is the image under $\iota$ of the
  point of $(q, z')$ (\cref{def:pencil-d-configuration}), as $\iota$
  inserts its zeros before the last coordinate. So the hinge
  $C(b_u, b_v)$ at each edge of the panel-hinge framework of $(q, z)$ is
  the image under $\varphi$ of that of $(q, z')$
  (\cref{lem:pencil-d-embedding-join}), and
  \cref{lem:screw-embedding-rank} applies. The lifting space $L(q)$ is a
  subspace, so it contains the zero height.
\end{proof}

\begin{lemma}[Attaining in a lower dimension]
  \label{lem:pencil-d-x0-attains-of-le}
  \lean{Graph.X0AttainsGrade.of_le}
  \leanok
  \uses{def:pencil-d-x0-attains, def:D-deficiency}
  Let $K$ be infinite, let $G$ be connected, and let $2 \le d' \le d$
  with $\operatorname{def}_{d'}(G) = \operatorname{def}_d(G)$. If the
  general configuration of $G$ attains in dimension $d'$
  (\cref{def:pencil-d-x0-attains}), it attains in dimension $d$.
\end{lemma}
\begin{proof}
  \leanok
  \uses{lem:pencil-d-embed-config-rank, lem:pencil-d-x0-one-witness,
        lem:pencil-x0-main-picture-open}
  Let $P$ be the polynomial of \cref{def:pencil-d-x0-attains} in
  dimension $d'$. As $K$ is infinite, $P$ has a non-root, which is an
  admissible picture. So $G$ has no loop and each closed neighbourhood
  has at least three members, and some nonzero polynomial has only main
  pictures as non-roots (\cref{lem:pencil-x0-main-picture-open}). Its
  product with $P$ has a non-root $q_0$, a main picture with
  $P(q_0) \neq 0$. Some $z' \in L(q_0)^{d'-2}$ gives $(q_0, z')$ the rank
  $D'(|V(G)| - 1) - \operatorname{def}_{d'}(G)$. Appending zero heights
  gives $z \in L(q_0)^{d-2}$, and $(q_0, z)$ has rank
  $D(|V(G)| - 1) - \operatorname{def}_{d'}(G)
  = D(|V(G)| - 1) - \operatorname{def}_d(G)$
  (\cref{lem:pencil-d-embed-config-rank}). As $G$ has a body,
  \cref{lem:pencil-d-x0-one-witness} applies.
\end{proof}

\begin{lemma}[The general configuration attaining gives an adjacent-distinct realization, in every dimension]
  \label{lem:pencil-d-x0-attains-distinct}
  \lean{Graph.X0AttainsGrade.hasGradedDistinctPencilRealization}
  \leanok
  \uses{def:pencil-d-x0-attains, def:pencil-d-distinct-motive}
  Let $K$ be infinite. If the general configuration of $G$ attains in
  dimension $d$ (\cref{def:pencil-d-x0-attains}), then $G$ has an
  adjacent-distinct pencil realization in dimension $d$
  (\cref{def:pencil-d-distinct-motive}).
\end{lemma}
\begin{proof}
  \leanok
  \uses{lem:pencil-d-config-realization}
  As $K$ is infinite, the nonzero polynomial $P$ of
  \cref{def:pencil-d-x0-attains} has a non-root $q$. This picture is
  admissible, and some $z \in L(q)^{d-2}$ gives $(q, z)$ the rank
  $D(|V(G)| - 1) - \operatorname{def}_d(G)$. Apply
  \cref{lem:pencil-d-config-realization} to $(q, z)$.
\end{proof}

\begin{theorem}[The pencil in every dimension when $\operatorname{def}_2(G) \le 3$]
  \label{thm:pencil-d-deficiency-two-le-three}
  \lean{Graph.x0AttainsGrade_of_deficiency_two_le_three,
        Graph.hasGradedDistinctPencilRealization_of_deficiency_two_le_three}
  \leanok
  \uses{def:pencil-x0-standing, def:pencil-d-x0-attains,
        def:pencil-d-distinct-motive, def:D-deficiency}
  Let $K$ be an infinite field and $d \ge 3$, and let $G$ satisfy the
  standing hypotheses (\cref{def:pencil-x0-standing}) with
  $\operatorname{def}_2(G) \le 3$. Then the general configuration of $G$
  attains in dimension $d$ (\cref{def:pencil-d-x0-attains}), and $G$ has
  an adjacent-distinct pencil realization in dimension $d$
  (\cref{def:pencil-d-distinct-motive}).
\end{theorem}
\begin{proof}
  \leanok
  \uses{thm:pencil-x0-coverage, lem:pencil-d-x0-attains-three,
        lem:deficiency-stable, lem:pencil-d-x0-attains-of-le,
        lem:pencil-d-x0-attains-distinct}
  The general configuration of $G$ attains in dimension three
  (\cref{thm:pencil-x0-coverage}, \cref{lem:pencil-d-x0-attains-three}).
  As $G$ is connected and
  $\operatorname{def}_2(G) + 3 \le 6 = \binom{4}{2}$,
  \cref{lem:deficiency-stable} with $d' = 3$ gives
  $\operatorname{def}_3(G) = \operatorname{def}_d(G)$. So it attains in dimension $d$
  (\cref{lem:pencil-d-x0-attains-of-le}), and
  \cref{lem:pencil-d-x0-attains-distinct} gives the realization.
\end{proof}

\subsection{The deficiency at every screw dimension}
\label{sec:pencil-d-deficiency}

The estimates of \cref{sec:main-component-sparse} that compare two
dimensions are stated there at $d = 3$. They hold in every dimension
$d$, with $\operatorname{def}_d$ in place of $\operatorname{def}_3$ and
$D - 1$ in place of $5$, where $D = \binom{d+1}{2}$. At $d = 3$ they are
\cref{lem:deficiency-core-bound}, part~(2) of
\cref{lem:deficiency-one-body-chain} and
\cref{lem:deficiency-two-body-chain}. Two facts replace the comparison
of the planar value with the value at $d = 3$. A partition of value at
most zero at $d = 2$ has value at most zero at every $d \ge 2$
(\cref{lem:deficiency-planar-nonpos}). A graph of deficiency zero in
some dimension has deficiency zero in every higher one
(\cref{lem:deficiency-rigid-higher}). We keep the planar notions of
\cref{sec:main-component-sparse}: the singleton value $s$, sparse graphs
and tight sets. A tight set $X$ has $\operatorname{def}_3(G[X]) = 0$
(\cref{lem:deficiency-tight-rigid}), so $\operatorname{def}_d(G[X]) = 0$
for every $d \ge 3$. We write $v_d(P)$ for the value of a partition $P$ at
$d$, and $d_G(P)$ for the number of edges of $G$ joining distinct parts
of $P$.

\begin{lemma}[A planar value at most zero is at most zero in every dimension]
  \label{lem:deficiency-planar-nonpos}
  \lean{Graph.three_mul_partitionDef_eq,
        Graph.partitionDef_nonpos_of_partitionDef_two_nonpos}
  \leanok
  \uses{def:D-deficiency}
  Let $d \ge 2$, let $D = \binom{d+1}{2}$, and let $P$ be a partition of
  $V(G)$. Then
  \[
    3\,v_d(P) \;=\; D\,v_2(P) - (D - 3)\,d_G(P).
  \]
  In particular, if $v_2(P) \le 0$, then $v_d(P) \le 0$.
\end{lemma}
\begin{proof}
  \leanok
  As $v_2(P) = 3(|P| - 1) - 2\,d_G(P)$, both sides equal
  $3D(|P| - 1) - 3(D - 1)\,d_G(P)$. The second claim follows, as
  $D \ge 3$ and $d_G(P) \ge 0$.
\end{proof}

\begin{lemma}[A graph of deficiency zero has deficiency zero in every higher dimension]
  \label{lem:deficiency-rigid-higher}
  \lean{Graph.deficiency_eq_zero_of_deficiency_eq_zero_of_le}
  \leanok
  \uses{def:D-deficiency}
  Let $1 \le d \le d'$, and let $G$ have a body. If
  $\operatorname{def}_d(G) = 0$, then $\operatorname{def}_{d'}(G) = 0$.
\end{lemma}
\begin{proof}
  \leanok
  \uses{lem:deficiency-zero-connected, lem:deficiency-antitone}
  At $d \ge 1$ we have $D \ge 1$, so $G$ is connected
  (\cref{lem:deficiency-zero-connected}). Hence
  $\operatorname{def}_{d'}(G) \le \operatorname{def}_d(G) = 0$
  (\cref{lem:deficiency-antitone}), and the deficiency of a graph with a
  body is nonnegative.
\end{proof}

\begin{lemma}[A maximal rigid set bounds the singleton values above it, in every dimension]
  \label{lem:deficiency-core-bound-d}
  \lean{Graph.deficiency_induce_eq_zero_of_le,
        Graph.partitionDef_induce_id_le_of_maximal_of_three_le_bodyBarDim}
  \leanok
  \uses{def:D-deficiency}
  Let $d \ge 2$, and let $G$ be sparse.
  \begin{enumerate}
    \item Let $W \subseteq Y \subseteq V(G)$ with $W$ nonempty and
      $\operatorname{def}_d(G[W]) = 0$. If $s(Y) \le s(W)$, and
      $s(W) \le s(Q)$ for every $Q$ with $W \subseteq Q \subsetneq Y$,
      then $\operatorname{def}_d(G[Y]) = 0$.
    \item Let $W \subseteq V(G)$ be nonempty with
      $\operatorname{def}_d(G[W]) = 0$, and let $X_0$ be empty or a
      single body outside $W$ that sends at most two edges into every
      set avoiding it and at most one into $W$. Suppose $W$ is maximal
      among the sets $W' \supseteq W$ in $V(G) \setminus X_0$ with
      $\operatorname{def}_d(G[W']) = 0$. Then $s(W) \le s(X)$ for every
      $X$ with $W \subseteq X \subseteq V(G)$.
  \end{enumerate}
\end{lemma}
\begin{proof}
  \leanok
  \uses{lem:deficiency-merge-rigid, lem:deficiency-singleton-bound,
        lem:deficiency-add-body, lem:deficiency-planar-nonpos}
  (1) Let $P$ be a partition of $Y$, with values taken in $G[Y]$. Let
  $P'$ be a partition coarser than $P$ with $W$ inside one part $Q$ and
  $v_d(P) \le v_d(P')$ (\cref{lem:deficiency-merge-rigid}). If $Q = Y$
  then $v_d(P') = 0$. Otherwise $v_2(P') \le s(Y) - s(Q) \le 0$
  (\cref{lem:deficiency-singleton-bound}), so $v_d(P') \le 0$
  (\cref{lem:deficiency-planar-nonpos}).

  (2) The proof of part~(2) of \cref{lem:deficiency-core-bound} holds
  with $\operatorname{def}_d$ in place of $\operatorname{def}_3$. It
  compares singleton values (\cref{lem:deficiency-add-body}), and it
  applies part~(1), which holds in dimension $d$.
\end{proof}

\begin{lemma}[The merged deficiency at a body of degree two, in every dimension]
  \label{lem:deficiency-one-body-chain-d}
  \lean{Graph.deficiencyMerged_add_bodyBarDim_sub_one_le}
  \leanok
  \uses{def:D-deficiency, def:deficiency-merged}
  Let $d \ge 2$ and $D = \binom{d+1}{2}$. Let $V_1 \subseteq V(G)$ and
  $x \in V(G) \setminus V_1$, and let $a \neq b$ in $V_1$ be joined to
  $x$ by edges. Suppose $\operatorname{def}_d(G[Y]) \neq 0$ for every
  set $Y$ of at least two bodies containing $x$. Then
  $\operatorname{def}_d(G[V_1]; a, b) + D - 1 \le
  \operatorname{def}_d(G[V_1])$.
\end{lemma}
\begin{proof}
  \leanok
  \uses{lem:deficiency-add-body, lem:separated-partition-value}
  Suppose otherwise, and let $P$ attain the merged deficiency $m$, with
  $Q$ the part of $a$ and $b$. Refining $Q$ by a partition $R$ of $Q$
  changes the value of $P$ by the value of $R$ in $G[Q]$
  (\cref{lem:separated-partition-value}). Since
  $\operatorname{def}_d(G[V_1]) \le m + D - 2$, every such $R$ has value
  at most $D - 2$ in $G[Q]$, and at most $0$ if it keeps $a$ and $b$
  together.

  We show that $\operatorname{def}_d(G[Q \cup \{x\}]) = 0$, against the
  hypothesis. Let $R'$ be a partition of $Q \cup \{x\}$, and compare its
  value with that of its restriction to $Q$
  (\cref{lem:deficiency-add-body}). If $x$ is alone, its part adds $D$
  and its two edges to $a$ and $b$ cost $2(D - 1)$, so the value is at
  most $(D - 2) + D - 2(D - 1) = 0$. If $x$ shares a part, each edge
  from $x$ to another part costs $D - 1$. When $a$ and $b$ are apart, at
  least one such edge exists, and the value is at most
  $(D - 2) - (D - 1) < 0$. When they are together, the value is at
  most $0$.
\end{proof}

\begin{lemma}[A non-adjacent pair outside every rigid set, in every dimension]
  \label{lem:deficiency-two-body-chain-d}
  \lean{Graph.deficiencyMerged_two_add_two_le_of_three_le}
  \leanok
  \uses{def:D-deficiency, def:deficiency-merged}
  Let $d \ge 3$, let $G$ be sparse, and let $a \neq b$ be bodies of $G$
  that are not adjacent, with $\operatorname{def}_d(G; a, b) \neq
  \operatorname{def}_d(G)$. Then $\operatorname{def}_2(G; a, b) + 2 \le
  \operatorname{def}_2(G)$.
\end{lemma}
\begin{proof}
  \leanok
  \uses{lem:deficiency-singleton-bound, lem:deficiency-add-body,
        lem:deficiency-tight-rigid, lem:deficiency-rigid-higher,
        lem:deficiency-merge-rigid}
  Otherwise, as in the proof of \cref{lem:deficiency-two-body-chain},
  the part $Z$ of $a$ and $b$ in a partition attaining the planar merged
  deficiency is a tight set, so $\operatorname{def}_3(G[Z]) = 0$
  (\cref{lem:deficiency-tight-rigid}). As $d \ge 3$,
  $\operatorname{def}_d(G[Z]) = 0$ (\cref{lem:deficiency-rigid-higher}).
  Then $\operatorname{def}_d(G; a, b) = \operatorname{def}_d(G)$
  (\cref{lem:deficiency-merge-rigid}).
\end{proof}

\begin{lemma}[A part of a tight partition has deficiency zero]
  \label{lem:deficiency-tight-part}
  \lean{Graph.IsTightPartition.deficiency_induce_eq_zero}
  \leanok
  \uses{def:tight-partition}
  Let $P$ be a partition of $V(G)$ of value
  $\operatorname{def}_d(G)$ (\cref{def:tight-partition}). Then
  $\operatorname{def}_d(G[Z]) = 0$ for every part $Z$ of $P$.
\end{lemma}
\begin{proof}
  \leanok
  \uses{lem:separated-partition-value}
  If $Z = V(G)$, then $P$ has the one part $V(G)$ and value $0$, so
  $\operatorname{def}_d(G) = 0$. Otherwise let $R$ be a partition of
  $Z$, and let $P'$ be $P$ with $Z$ replaced by the parts of $R$. Both
  $P$ and $P'$ have each part inside $Z$ or inside $V(G) \setminus Z$,
  and they agree outside $Z$. So their values differ by the values of
  $R$ and of $\{Z\}$ in $G[Z]$ (\cref{lem:separated-partition-value}),
  and the second is zero. As $P$ attains the deficiency, the value of
  $R$ in $G[Z]$ is at most zero, and $\{Z\}$ attains it.
\end{proof}

With these estimates the existence of a core when no chain is usable
(\cref{thm:pencil-x0-theorem-s}) holds in every dimension, with
rigidity read at $D$ and additivity in the plane. Let $d \ge 2$. A set
$W$ of at least two bodies is \emph{rigid in dimension $d$} if
$\operatorname{def}_d(G[W]) = 0$. A \emph{core in dimension $d$} of $G$
is a set $W \subsetneq V(G)$ of at least two bodies, rigid in dimension
$d$, such that no body outside $W$ is adjacent to two bodies of $W$,
together with a body $r \in W$ such that
\[
  \operatorname{def}_2(G[W]) + \operatorname{def}_2(G/G[W]) \;\le\;
  \operatorname{def}_2(G),
\]
where $G/G[W]$ contracts $W$ to $r$ (\cref{def:rigid-contraction}). At
$d = 3$ it is the core of \cref{thm:pencil-x0-theorem-s}. The chains
usable in dimension $d$ read the merged deficiencies at $D$, with
$D - 1$ in place of $5$, and keep the planar condition at two interior
bodies.

\begin{definition}[Usable chains in every dimension]
  \label{def:pencil-d-usable-chain}
  \lean{Graph.ChainUsableAt}
  \leanok
  \uses{def:pencil-x0-chain, def:deficiency-merged}
  Let $d \ge 2$ and $D = \binom{d+1}{2}$. A chain
  $a - x_1 - \cdots - x_k - b$ on $V_1$ (\cref{def:pencil-x0-chain}) is
  \emph{usable in dimension $d$} if $k \ge 3$; or $k = 2$ and either
  $\operatorname{def}_d(G[V_1]; a, b) = \operatorname{def}_d(G[V_1])$, or
  $a$ and $b$ are not adjacent and $\operatorname{def}_2(G[V_1]; a, b) +
  2 \le \operatorname{def}_2(G[V_1])$; or $k = 1$ and either
  $\operatorname{def}_d(G[V_1]; a, b) = \operatorname{def}_d(G[V_1])$ or
  $\operatorname{def}_d(G[V_1]; a, b) + D - 1 \le
  \operatorname{def}_d(G[V_1])$. At $d = 3$ these are the usable chains
  of \cref{def:pencil-x0-usable-chain}.
\end{definition}

\begin{lemma}[A maximal rigid set is a core, in every dimension]
  \label{lem:pencil-d-maximal-rigid-core}
  \lean{Graph.IsX0Graph.exists_additiveCore_of_rigidAt}
  \leanok
  \uses{def:pencil-x0-standing, def:rigid-contraction}
  Let $d \ge 2$, and let $G$ satisfy the standing hypotheses and be
  sparse. Let $X_0$ be empty or a single body of degree at most two that
  has at most one neighbour in every set rigid in dimension $d$ avoiding
  it, and let $X_0$ be nonempty if $\operatorname{def}_d(G) = 0$. Let
  $Y$ be a set of at least two bodies, rigid in dimension $d$ and
  disjoint from $X_0$. Then $G$ has a core in dimension $d$ that avoids
  $X_0$.
\end{lemma}
\begin{proof}
  \leanok
  \uses{lem:deficiency-add-body, lem:deficiency-core-bound-d,
        lem:deficiency-additive-core}
  As in the proof of \cref{thm:pencil-x0-theorem-s}, let $W$ be maximal
  among the sets rigid in dimension $d$ that contain $Y$ and avoid
  $X_0$, and let $r \in W$. Then $W \neq V(G)$, as $X_0$ is nonempty
  when $G$ is rigid in dimension $d$. A body outside
  $W$ and $X_0$ with two neighbours in $W$ would enlarge $W$
  (\cref{lem:deficiency-add-body}, at $D \ge 3$), and the body of $X_0$
  has at most one. By part~(2) of \cref{lem:deficiency-core-bound-d},
  $s(W) \le s(X)$ for every $X \supseteq W$, and
  \cref{lem:deficiency-additive-core} gives the inequality.
\end{proof}

\begin{theorem}[A core when no chain is usable, in every dimension]
  \label{thm:pencil-d-theorem-s}
  \lean{Graph.IsX0Graph.exists_additiveCoreAt}
  \leanok
  \uses{def:pencil-d-usable-chain, def:pencil-x0-standing,
        def:rigid-contraction}
  Let $d \ge 3$, and let $G$ satisfy the standing hypotheses, with
  $G - v$ connected for every body $v$, some body of degree other than
  two, and no planar-rigid set. If no chain of $G$ is usable in
  dimension $d$, then $G$ has a core in dimension $d$.
\end{theorem}
\begin{proof}
  \leanok
  \uses{lem:pencil-d-maximal-rigid-core, lem:deficiency-sparse,
        lem:pencil-x0-sparse-count, lem:pencil-x0-chain-exists,
        lem:deficiency-one-body-chain-d, lem:deficiency-two-body-chain-d,
        lem:deficiency-tight-rigid, lem:deficiency-rigid-higher,
        lem:deficiency-merge-rigid, lem:deficiency-zero-connected}
  The proof of \cref{thm:pencil-x0-theorem-s}, with rigidity in
  dimension $d$. A set rigid in dimension three is rigid in dimension $d$
  (\cref{lem:deficiency-rigid-higher}). As $G$ is sparse
  (\cref{lem:deficiency-sparse}), it has bodies of degree two, each on a
  chain (\cref{lem:pencil-x0-sparse-count},
  \cref{lem:pencil-x0-chain-exists}). As no chain is usable, each has at
  most two interior bodies. A chain $a - x_1 - x_2 - b$ on
  $V_1$ has adjacent ends: otherwise
  $\operatorname{def}_d(G[V_1]; a, b) \neq \operatorname{def}_d(G[V_1])$,
  and \cref{lem:deficiency-two-body-chain-d} would make it usable. Its
  $4$-cycle is tight, so rigid in dimension three
  (\cref{lem:deficiency-tight-rigid}), and so in dimension $d$.

  If $\operatorname{def}_d(G) > 0$, a chain gives a set rigid in
  dimension $d$: its $4$-cycle, or, at one interior body $x$, a set
  containing $x$, as otherwise the chain would be usable
  (\cref{lem:deficiency-one-body-chain-d}). Then
  \cref{lem:pencil-d-maximal-rigid-core} with $X_0$ empty gives the core.

  If $\operatorname{def}_d(G) = 0$, take a body $x$ of degree two and a
  set rigid in dimension $d$ that avoids it: off the $4$-cycle of a chain
  with two interior bodies if two bodies of degree two are adjacent
  (\cref{lem:pencil-x0-sparse-count} gives four bodies of degree two),
  and otherwise inside $V(G) \setminus \{x\}$, which contains a set rigid
  in dimension three as $s_6(V(G) \setminus \{x\}) \le 0$
  (\cref{lem:pencil-x0-sparse-count}, \cref{lem:deficiency-sparse}). A
  neighbour of $x$ in a set $W$ rigid in dimension $d$ avoiding $x$ has
  two neighbours in $W$ (\cref{lem:deficiency-zero-connected}), so degree
  at least three. Two such neighbours $a$ and $b$ would make $a - x - b$
  a chain on $V(G) \setminus \{x\}$, usable as
  $\operatorname{def}_d(G - x; a, b) = \operatorname{def}_d(G - x)$
  (\cref{lem:deficiency-merge-rigid}). So
  \cref{lem:pencil-d-maximal-rigid-core} with $X_0 = \{x\}$ gives the
  core.
\end{proof}

\begin{lemma}[A core when the deficiency in dimension three is positive]
  \label{lem:pencil-d-core-positive-deficiency}
  \lean{Graph.IsX0Graph.exists_core_of_deficiency_three_pos}
  \leanok
  \uses{def:pencil-x0-standing, def:rigid-contraction}
  Let $d \ge 3$, and let $G$ satisfy the standing hypotheses and be
  sparse, with no two bodies of degree two adjacent. If
  $\operatorname{def}_3(G) > 0$, then $G$ has a core in dimension $d$.
\end{lemma}
\begin{proof}
  \leanok
  \uses{lem:deficiency-tight-part, lem:deficiency-rigid-higher,
        lem:pencil-x0-sparse-count, lem:deficiency-planar-nonpos,
        lem:deficiency-singleton-bound,
        lem:deficiency-core-bound-d, lem:deficiency-add-body,
        lem:deficiency-additive-core}
  Let $P$ be a partition of value $\operatorname{def}_3(G) > 0$. Each of
  its parts is rigid in dimension three
  (\cref{lem:deficiency-tight-part}), so in dimension $d$
  (\cref{lem:deficiency-rigid-higher}). The partition into singletons
  has value $6(|V(G)| - 1) - 5|E(G)| < 0$, as $5|E(G)| \ge 6|V(G)|$ when
  no two bodies of degree two are adjacent (the count in the proof of
  \cref{lem:pencil-x0-sparse-count}). So $P$ has a part $Z$ of at least
  two bodies, and $Z$ is rigid in dimension $d$. As
  $v_3(P) = 2\,v_2(P) - d_G(P) > 0$
  (\cref{lem:deficiency-planar-nonpos}), $P$ has planar value at least
  one, so $s(Z) \le s(V(G)) - 1$ (\cref{lem:deficiency-singleton-bound}).
  Let $W$ have least singleton value $s(W)$ among the sets rigid in
  dimension $d$, and let $r \in W$. Then $s(W) \le s(Z) < s(V(G))$, and
  $W \neq V(G)$.

  Suppose that some $X \supseteq W$ has $s(X) < s(W)$, and take one of
  least size. Then $s(W) \le s(Q)$ whenever $W \subseteq Q \subsetneq X$,
  so $X$ is rigid in dimension $d$ (part~(1) of
  \cref{lem:deficiency-core-bound-d}), against the choice of $W$. So
  $s(W) \le s(X)$ for every $X \supseteq W$. A body outside $W$ with two
  neighbours in it would lower $s$ on joining $W$
  (\cref{lem:deficiency-add-body}), and
  \cref{lem:deficiency-additive-core} gives the inequality.
\end{proof}

\begin{lemma}[A core at an open ear of two or more interior bodies]
  \label{lem:pencil-d-chain-core}
  \lean{Graph.IsOpenEar.exists_additiveCoreAt_of_partitionDef_nonpos}
  \leanok
  \uses{def:open-ear, def:pencil-x0-standing, def:rigid-contraction}
  Let $d \ge 2$ and $D = \binom{d+1}{2}$, and let $G$ satisfy the
  standing hypotheses and be sparse. Let
  $a - x_1 - \cdots - x_k - b$ be an open ear of $G$ on $V_1$
  (\cref{def:open-ear}) with $k \ge 2$ and $s_D(V_1) \le 0$. Then $G$
  has a core in dimension $d$ that avoids $x_1$.
\end{lemma}
\begin{proof}
  \leanok
  \uses{lem:deficiency-sparse, lem:deficiency-zero-connected,
        lem:pencil-d-maximal-rigid-core}
  As $a \neq b$ lie in $V_1$ and $s_D(V_1) \le 0$, $V_1$ contains a set
  $Y$ rigid in dimension $d$ (\cref{lem:deficiency-sparse}), which
  avoids $x_1$. As $G$ is simple, the interior bodies have degree two,
  and the neighbours of $x_1$ are $a$ and $x_2$. A set $W$ rigid in
  dimension $d$ that avoids $x_1$ does not contain $x_2$, since $x_2$
  would have two neighbours in $W$ (\cref{lem:deficiency-zero-connected})
  and one of its two neighbours is $x_1$. So $x_1$ has at most one
  neighbour in $W$, and \cref{lem:pencil-d-maximal-rigid-core} with
  $X_0 = \{x_1\}$ gives the core.
\end{proof}

\subsection{The induction in every dimension}
\label{sec:pencil-d-induction}

From here on we extend the induction of \cref{sec:main-component} to
every dimension $d \ge 3$, to show that the general configuration of
every graph satisfying the standing hypotheses
(\cref{def:pencil-x0-standing}) attains in dimension $d$
(\cref{def:pencil-d-x0-attains}). At $d = 3$ this is
\cref{thm:pencil-x0-coverage}. As there, we argue by strong induction on
the number of bodies, at a fixed $d$, with all $d - 2$ heights of a
configuration general. A graph also attains in dimension $d$ when it
attains in a lower dimension $d'$ with
$\operatorname{def}_{d'}(G) = \operatorname{def}_d(G)$
(\cref{lem:pencil-d-x0-attains-of-le}).

The planar objects of \cref{sec:main-component-carrier} serve in every
dimension unchanged: the pictures and their admissibility, the lifting
space $L(q)$ and its lifting system, and the main pictures
(\cref{def:pencil-main-picture}, \cref{lem:pencil-x0-main-picture-open}).
So does Jackson and Jord\'an's equality (\cref{thm:pencil-jj-equality}),
which concerns $L(q)$ alone. A configuration in dimension $d$ takes each
of its $d - 2$ heights from $L(q)$. The objects that read the points of
a configuration are restated in \cref{sec:pencil-d-carrier}. This
section fixes the notation, and shows that over an admissible picture
the pencil configurations are those with every height in $L(q)$
(\cref{lem:pencil-d-condition-linear}). The next section computes the
rank of a flat configuration, which settles the graphs whose two
deficiencies agree (\cref{sec:pencil-d-flat}). The steps of the
induction follow, beginning with a cut vertex and a chain of bridges
(\cref{sec:pencil-d-cut}).

We fix the following notation and terms, for $d \ge 3$ and
$D = \binom{d+1}{2}$. At $d = 3$ they are those of
\cref{sec:main-component}.
\begin{itemize}
  \item $\operatorname{def}_d(G)$ is the deficiency at $D$
    (\cref{def:D-deficiency}), and $\operatorname{def}_2(G)$, the
    deficiency at $3$, is the \emph{planar deficiency}. On a connected
    graph $\operatorname{def}_d(G) \le \operatorname{def}_2(G)$
    (\cref{lem:deficiency-antitone}).
  \item The \emph{target} of $G$ in dimension $d$ is its deficiency rank
    $\operatorname{tgt}_d(G) = D(|V(G)| - 1) - \operatorname{def}_d(G)$
    (\cref{def:pencil-d-rank-hypothesis}). We say that $G$ \emph{has its
    target} at a configuration in dimension $d$ of rank
    $\operatorname{tgt}_d(G)$ (\cref{def:pencil-d-configuration}).
  \item $L(q)^{d-2}$ is the space of tuples $z = (z^1, \dots, z^{d-2})$
    of heights in $L(q)$ (\cref{def:pencil-lifting-space}), a linear
    subspace of the space of all tuples of $d - 2$ heights. A
    configuration $(q, z)$ with $q$ admissible and $z \in L(q)^{d-2}$ is
    \emph{attaining} if $G$ has its target at it, and its heights $z$
    are then \emph{attaining heights}.
  \item The configurations in dimension $d$ at which the points of every
    closed neighbourhood span a subspace of dimension at most three form
    the \emph{pencil configuration space} of $G$. Over an admissible
    picture $q$ they are the configurations $(q, z)$ with
    $z \in L(q)^{d-2}$ (\cref{lem:pencil-d-condition-linear}). The
    \emph{main component} is the closure of the part of this space over
    the main pictures. As at $d = 3$, we never form it
    (\cref{rem:pencil-x0-main-component}).
  \item \emph{General} means off the zero set of a nonzero polynomial: in
    the picture coordinates, and over each picture in the heights
    $z \in L(q)^{d-2}$. Over an infinite field, two polynomials that are
    each nonzero somewhere on $L(q)^{d-2}$ are nonzero together somewhere
    on it (\cref{lem:pencil-x0-fibre-intersection}).
  \item A \emph{flat configuration} is a configuration $(q, 0)$ whose
    $d - 2$ heights are all zero.
  \item A set $W$ of at least two bodies is \emph{rigid in dimension $d$}
    if $\operatorname{def}_d(G[W]) = 0$, and \emph{planar-rigid} if
    $\operatorname{def}_2(G[W]) = 0$. A set rigid in some dimension is
    rigid in every higher one (\cref{lem:deficiency-rigid-higher}). In
    particular a planar-rigid set is rigid in dimension $d$.
  \item A contraction step contracts a set $W \subsetneq V(G)$ of at
    least two bodies, the \emph{core}, to one of its bodies. We write
    $H = G[W]$ and $G/H$ for the contraction
    (\cref{def:rigid-contraction}). The core is \emph{additive} in
    dimension $d$ if it is rigid in dimension $d$ and
    $\operatorname{def}_2(H) + \operatorname{def}_2(G/H) \le
    \operatorname{def}_2(G)$. So rigidity is read at $D$, and additivity
    in the plane.
  \item A graph is \emph{smaller than} $G$ if it satisfies the standing
    hypotheses and has fewer bodies than $G$.
\end{itemize}

At $d = 3$ the points of a closed neighbourhood are coplanar exactly
when the height agrees with an affine function of the picture there
(\cref{lem:pencil-condition-linear}). In dimension $d$ the same holds
for each of the $d - 2$ heights.

\begin{lemma}[Coplanarity is linear in the heights, in every dimension]
  \label{lem:pencil-d-condition-linear}
  \uses{def:pencil-d-configuration, def:pencil-lifting-space,
        def:pencil-admissible-picture}
  Let $d \ge 3$, let $q$ be an admissible picture of $G$ and
  $z^1, \dots, z^{d-2}$ heights, and let $v$ be a body of $G$. The points
  of the configuration $(q, z)$ (\cref{def:pencil-d-configuration}) at
  the members of the closed neighbourhood of $v$ span a subspace of
  $K^{d+1}$ of dimension at most three if and only if each $z^i$ agrees
  on that neighbourhood with an affine function
  $a^i x + c^i y + e^i$ of the picture. They then span the
  three-dimensional image $\Pi_v$ of the injective linear map
  \[
    \sigma_v(\xi, \eta, \omega) = (\xi,\ \eta,\
    a^1 \xi + c^1 \eta + e^1 \omega,\ \dots,\
    a^{d-2} \xi + c^{d-2} \eta + e^{d-2} \omega,\ \omega),
  \]
  and $\Pi_v^{\perp}$ is spanned by the $d - 2$ vectors $n^i$ with
  $n^i \cdot (\xi, \eta, \zeta^1, \dots, \zeta^{d-2}, \omega) =
  a^i \xi + c^i \eta - \zeta^i + e^i \omega$. In particular the points of
  every closed neighbourhood span a subspace of dimension at most three
  if and only if every $z^i$ lies in $L(q)$
  (\cref{def:pencil-lifting-space}).
\end{lemma}
\begin{proof}
  Write $b$ for the points of $(q, z)$, and let
  $\pi \colon K^{d+1} \to K^3$ keep the first, the second and the last
  coordinate. Then $\pi(b_w) = (x_w, y_w, 1)$ is the homogeneous picture
  point of $w$. As $q$ is admissible, three members of the closed
  neighbourhood of $v$ have linearly independent homogeneous picture
  points. Their points are therefore linearly independent, and the points
  of the closed neighbourhood span a subspace $S$ with $\pi(S) = K^3$.

  Suppose that $\dim S \le 3$. Then $\dim S = 3$, and $\pi$ restricts to
  an isomorphism of $S$ onto $K^3$. Its inverse is a linear map
  $\sigma \colon K^3 \to K^{d+1}$ with $\pi \circ \sigma$ the identity.
  So $\sigma$ keeps the first, the second and the last coordinate, and its
  coordinate of the $i$-th height is a linear form
  $a^i \xi + c^i \eta + e^i \omega$. Each point $b_w$ of the closed
  neighbourhood lies in $S$, and $\pi(b_w) = (x_w, y_w, 1)$, so
  $b_w = \sigma(x_w, y_w, 1)$. Its $i$-th height is
  $z^i_w = a^i x_w + c^i y_w + e^i$.

  Conversely, suppose that $z^i_w = a^i x_w + c^i y_w + e^i$ for every
  $w$ in the closed neighbourhood and every $i$. Then
  $b_w = \sigma_v(x_w, y_w, 1)$ there. The map $\sigma_v$ is injective,
  as $\pi \circ \sigma_v$ is the identity, so its image $\Pi_v$ has
  dimension three. It contains $S$, and $S = \Pi_v$ since
  $\pi(S) = K^3$. Each $n^i$ is orthogonal to $\Pi_v$, as
  \[
    n^i \cdot \sigma_v(\xi, \eta, \omega) = a^i \xi + c^i \eta
    - (a^i \xi + c^i \eta + e^i \omega) + e^i \omega = 0.
  \]
  The vectors $n^i$ are linearly independent, since $n^i$ alone has a
  nonzero coordinate at the $i$-th height. As the dot product of
  $K^{d+1}$ is nondegenerate, $\Pi_v^{\perp}$ has dimension $d - 2$, so
  the $n^i$ span it. The last claim is the first at every body
  (\cref{def:pencil-lifting-space}).

  At $d = 3$ this is \cref{lem:pencil-condition-linear}. There the
  dimension enters in two places: the plane through the closed
  neighbourhood is given by one normal vector, and the height is found by
  solving the equation of that plane, which is not vertical. In
  dimension $d$ the span $S$ has codimension $d - 2$, and no single
  normal vector. The projection $\pi$ takes the place of the normal:
  that $\pi$ is injective on $S$ says that $S$ is not vertical, and the
  inverse of $\pi$ on $S$ gives all $d - 2$ heights at once.
\end{proof}

\subsection{The flat rank in every dimension}
\label{sec:pencil-d-flat}

This section computes the rank of a flat configuration $(q, 0)$ in
dimension $d$ (\cref{thm:pencil-d-flat-rank}). As at $d = 3$
(\cref{sec:main-component-flat}), with Jackson and Jord\'an's equality
it settles the graphs whose two deficiencies agree
(\cref{cor:pencil-d-jj-flat}). The points of a flat configuration have
zero heights, and they are the homogeneous picture points with $d - 2$
zero coordinates inserted before the last. So the flat configuration in
dimension $d$ is the configuration of the picture in dimension two,
embedded by $\iota$ (\cref{def:pencil-d-embedding}), and its rank is
read off the plane framework of \cref{lem:pencil-lifting-planes-motions}.

Over an infinite field and under the standing hypotheses, Jackson and
Jord\'an's equality gives $\dim L(q) = 3 + \operatorname{def}_2(G)$ at a
general picture (\cref{thm:pencil-jj-equality}). The flat configuration
over such a picture then has rank
$D(|V(G)| - 1) - \operatorname{def}_2(G)$, which is the target exactly
when $\operatorname{def}_2(G) = \operatorname{def}_d(G)$. The later steps
use the equality at the graphs that they compare with $G$, through
$\dim L(q)$ alone.

\begin{theorem}[The flat rank, in every dimension]
  \label{thm:pencil-d-flat-rank}
  \uses{def:pencil-d-configuration, def:pencil-lifting-space,
        def:D-deficiency}
  Let $d \ge 3$ and $D = \binom{d+1}{2}$, let $G$ be connected, and let
  $q$ be an admissible picture of $G$. Then the flat configuration
  $(q, 0)$ in dimension $d$, all of whose $d - 2$ heights are zero
  (\cref{def:pencil-d-configuration}), has rank exactly
  \[
    D(|V(G)| - 1) + 3 - \dim L(q).
  \]
  Hence its rank is at most $D(|V(G)| - 1) - \operatorname{def}_2(G)$,
  with equality if and only if $\dim L(q) = 3 + \operatorname{def}_2(G)$.
\end{theorem}
\begin{proof}
  \uses{lem:pencil-d-embed-config-rank, lem:pencil-lifting-planes-motions,
        lem:pencil-lifting-space-deficiency}
  The configuration of $q$ in dimension two has no heights, and its point
  at each body $v$ is the homogeneous picture point $(x_v, y_v, 1)$
  (\cref{def:pencil-d-configuration}). So its rank is the rank of the
  plane framework $P_q$ with these panel normals, which is
  $3|V(G)| - \dim L(q)$ (\cref{lem:pencil-lifting-planes-motions}). The
  flat configuration in dimension $d$ appends $d - 2$ zero heights to it.
  As $G$ is connected, this adds $(D - 3)(|V(G)| - 1)$ to the rank
  (\cref{lem:pencil-d-embed-config-rank}, with $d' = 2$). The sum is
  \[
    3|V(G)| - \dim L(q) + (D - 3)(|V(G)| - 1)
    = D(|V(G)| - 1) + 3 - \dim L(q).
  \]
  The bound and its equality case follow from
  $\dim L(q) \ge 3 + \operatorname{def}_2(G)$
  (\cref{lem:pencil-lifting-space-deficiency}).

  At $d = 3$ this is \cref{thm:pencil-flat-rank}, and its proof gives the
  same count in every dimension. In dimension $d$ the hinge at an edge
  $uv$ is the image of the join $b_u \vee b_v$ under the complement
  isomorphism (\cref{def:panel-support-extensor}). Its inverse takes the
  motions to the assignments of elements of $\bigwedge^2 K^{d+1}$ to the
  bodies whose difference along each edge $uv$ is a multiple of
  $b_u \vee b_v$. Let $\hat H \subseteq K^{d+1}$ be the
  three-dimensional subspace of vectors with zero heights. The joins of
  the points of $(q, 0)$ lie in $\bigwedge^2 \hat H$. As at $d = 3$, the
  component of such an assignment in a complement of
  $\bigwedge^2 \hat H$ is constant, as $G$ is connected, and its
  component in $\bigwedge^2 \hat H$ is a family of lifting planes
  (\cref{def:pencil-lifting-planes}). The complement has dimension
  $D - 3$, where it has dimension $3$ at $d = 3$. So the motions have
  dimension $D - 3 + \dim L(q)$, and rank--nullity in the
  $D|V(G)|$-dimensional space of screw assignments gives the same rank.
\end{proof}

\begin{corollary}[When the flat configuration has the target, in every dimension]
  \label{cor:pencil-d-flat-attains}
  \uses{def:pencil-d-configuration, def:pencil-d-x0-attains,
        def:D-deficiency}
  Let $d \ge 3$ and $D = \binom{d+1}{2}$, let $G$ be connected, and let
  $q$ be an admissible picture of $G$. The flat configuration $(q, 0)$ in
  dimension $d$, all of whose heights are zero, has rank
  $D(|V(G)| - 1) - \operatorname{def}_d(G)$ if
  and only if $\dim L(q) = 3 + \operatorname{def}_d(G)$, and this holds
  if and only if equality holds in
  \cref{lem:pencil-lifting-space-deficiency} at $q$ and
  $\operatorname{def}_2(G) = \operatorname{def}_d(G)$. If
  $\dim L(q) \le 3 + \operatorname{def}_d(G)$, then the general
  configuration of $G$ attains in dimension $d$
  (\cref{def:pencil-d-x0-attains}).
\end{corollary}
\begin{proof}
  \uses{thm:pencil-d-flat-rank, lem:pencil-lifting-space-deficiency,
        lem:deficiency-antitone, def:pencil-main-picture,
        lem:pencil-d-x0-one-witness}
  The first equivalence is the exact rank of
  \cref{thm:pencil-d-flat-rank}. As $G$ is connected and $2 \le d$, every
  admissible picture $q'$ has
  \[
    \dim L(q') \;\ge\; 3 + \operatorname{def}_2(G) \;\ge\;
    3 + \operatorname{def}_d(G)
  \]
  (\cref{lem:pencil-lifting-space-deficiency},
  \cref{lem:deficiency-antitone}). At $q' = q$ this gives the second
  equivalence.

  Now suppose that $\dim L(q) \le 3 + \operatorname{def}_d(G)$. Then
  $\dim L(q') \ge \dim L(q)$ at every admissible picture $q'$, so $q$ is
  a main picture (\cref{def:pencil-main-picture}). By the exact rank, the
  flat configuration $(q, 0)$ has rank at least
  $D(|V(G)| - 1) - \operatorname{def}_d(G)$, and its heights, all zero,
  lie in $L(q)^{d-2}$. As $G$ is connected, it has a body. So
  \cref{lem:pencil-d-x0-one-witness} applies, and the general
  configuration of $G$ attains in dimension $d$.

  At $d = 3$ this is \cref{cor:pencil-flat-attains}. There the dimension
  enters only through the flat rank and the comparison of
  $\operatorname{def}_2$ with $\operatorname{def}_3$. Both hold in
  dimension $d$, with $\operatorname{def}_d$ in place of
  $\operatorname{def}_3$.
\end{proof}

\begin{corollary}[The general configuration attains when the two deficiencies agree, in every dimension]
  \label{cor:pencil-d-jj-flat}
  \uses{def:pencil-d-x0-attains, def:D-deficiency}
  Let $K$ be infinite and $d \ge 3$, and let $G$ be connected and simple,
  with at least three members in every closed neighbourhood. If
  $\operatorname{def}_2(G) = \operatorname{def}_d(G)$, then the general
  configuration of $G$ attains in dimension $d$
  (\cref{def:pencil-d-x0-attains}).
\end{corollary}
\begin{proof}
  \uses{thm:pencil-jj-equality, cor:pencil-d-flat-attains}
  As $G$ is connected, it has a body. By Jackson and Jord\'an's equality
  (\cref{thm:pencil-jj-equality}), some main picture $q$ of $G$ has
  \[
    \dim L(q) = 3 + \operatorname{def}_2(G) = 3 + \operatorname{def}_d(G).
  \]
  So the general configuration of $G$ attains in dimension $d$
  (\cref{cor:pencil-d-flat-attains}). At $d = 3$ this is
  \cref{cor:pencil-jj-flat}. The equality is planar, and the proof uses
  it as it stands.
\end{proof}

The hypothesis of \cref{cor:pencil-d-jj-flat} does not depend on the
dimension, so the corollary also follows from its case $d = 3$.

\begin{remark}[Equal deficiencies in one dimension and in all]
  \label{rem:pencil-d-flat-class}
  Let $G$ be connected, and suppose that
  $\operatorname{def}_2(G) = \operatorname{def}_d(G) = t$ for one
  $d \ge 3$. Then $\operatorname{def}_{d'}(G) = t$ for every $d' \ge 2$.
  Indeed, for a partition $P$ of $V(G)$ let
  $c(P) = d_G(P) - (|P| - 1)$, which is nonnegative as $G$ is connected
  (as in the proof of \cref{lem:deficiency-stable}). The value of $P$ at
  a screw dimension $D'$ is then $|P| - 1 - (D' - 1)\,c(P)$. Let $P$
  attain $\operatorname{def}_d(G)$, and let $D = \binom{d+1}{2} \ge 6$.
  Then $|P| - 1 = t + (D - 1)\,c(P)$, so the planar value of $P$ is
  $t + (D - 3)\,c(P)$. It is at most $\operatorname{def}_2(G) = t$, so
  $c(P) = 0$. Then $|P| - 1 = t$, and $P$ has value $t$ at every $D'$.
  So $\operatorname{def}_{d'}(G) \ge t$, and
  $\operatorname{def}_{d'}(G) \le \operatorname{def}_2(G) = t$ for
  $d' \ge 2$ (\cref{lem:deficiency-antitone}). In particular
  $\operatorname{def}_3(G) = \operatorname{def}_d(G)$. So under the
  hypotheses of \cref{cor:pencil-d-jj-flat} the general configuration of
  $G$ attains in dimension three (\cref{cor:pencil-jj-flat},
  \cref{lem:pencil-d-x0-attains-three}), and hence in dimension $d$
  (\cref{lem:pencil-d-x0-attains-of-le}).
\end{remark}

\subsection{Cut vertices and bridges in every dimension}
\label{sec:pencil-d-cut}

This section proves the first two steps of the induction in dimension
$d$: at a cut vertex (\cref{thm:pencil-d-cut}) and along a chain of
bridges (\cref{thm:pencil-d-bridge}). At $d = 3$ they are
\cref{thm:pencil-x0-cut} and \cref{thm:pencil-x0-bridge}, and we keep
the three parts of their proofs (\cref{sec:main-component-cut}). We take
a main picture $q$ of $G$ off the zero sets of the polynomials of the
attainment at the two sides. We find heights $z \in L_G(q)^{d-2}$ at
which both sides have their targets. A rank-addition lemma and a
deficiency count then give $(q, z)$ rank at least
$\operatorname{tgt}_d(G)$, and the general configuration of $G$ attains
in dimension $d$ (\cref{lem:pencil-d-x0-one-witness}). The first part is
planar. In the second, the restriction of heights to a side, and their
extension back to $G$, act on each of the $d - 2$ heights separately. In
the third, the rank-addition lemmas and the deficiency identities hold
at every screw dimension $D$, and a bridge adds $D - 1$ to the rank
where it adds $5$ at $d = 3$.

For a subgraph $G'$ of $G$ and a configuration $(q, z)$ of $G$ in
dimension $d$, the \emph{rank of $G'$ at $(q, z)$} is the rank of the
panel-hinge framework on $G'$ whose normal at each body of $G'$ is its
point in $(q, z)$ (\cref{def:pencil-d-configuration}). We say that $G'$
\emph{has its target} at $(q, z)$ when this rank is
$\operatorname{tgt}_d(G')$. The next lemma reads the rank of $G'$ at the
restriction of $z$ to $V(G')$.

\begin{lemma}[Heights restrict to a subgraph, in every dimension]
  \label{lem:pencil-d-lifting-restrict}
  \uses{def:pencil-d-configuration, def:pencil-lifting-space,
        def:panel-hinge-framework}
  Let $d \ge 3$, let $G' \le G$, let $q$ be a picture of $G$, and let
  $z = (z^1, \dots, z^{d-2})$ be heights of $G$. Write $z'$ for the tuple
  whose $i$-th height agrees with $z^i$ on $V(G')$ and vanishes
  elsewhere.
  \begin{enumerate}
    \item If $z \in L_G(q)^{d-2}$, then $z' \in L_{G'}(q)^{d-2}$.
    \item The points of the configurations $(q, z')$ and $(q, z)$ in
      dimension $d$ agree at the bodies of $G'$. So the rank of $G'$ at
      $(q, z')$ is its rank at $(q, z)$.
    \item A polynomial in the heights of $G'$, read at $z'$, agrees with
      a polynomial in the heights of $G$, read at $z$.
  \end{enumerate}
\end{lemma}
\begin{proof}
  \uses{lem:pencil-lifting-restrict, lem:pencil-rank-congr}
  At $d = 3$ this is \cref{lem:pencil-lifting-restrict}, and in dimension
  $d$ each claim holds height by height.

  (1) A tuple of heights lies in $L_G(q)^{d-2}$ exactly when each of its
  heights lies in $L_G(q)$. Each $z'^i$ is the restriction of
  $z^i \in L_G(q)$, so it lies in $L_{G'}(q)$ by part~(1) of
  \cref{lem:pencil-lifting-restrict}.

  (2) The point of $(q, z')$ at a body $w$ is
  $(x_w, y_w, z'^1_w, \dots, z'^{d-2}_w, 1)$
  (\cref{def:pencil-d-configuration}). It reads only the picture point and
  the heights at $w$, and at a body $w$ of $G'$ each $z'^i_w$ equals
  $z^i_w$. The rank of a panel-hinge framework on $G'$ depends only on its
  normals at the bodies of $G'$ (\cref{lem:pencil-rank-congr}). That lemma
  concerns panel-hinge frameworks in every dimension, and its proof does
  not use the dimension. So the two ranks agree.

  (3) Substituting zero for the heights $z^i_w$ with $w \notin V(G')$ in a
  polynomial and then evaluating at $z$ is the same as evaluating the
  polynomial at $z'$, which already vanishes there. This is the
  substitution of part~(3) of \cref{lem:pencil-lifting-restrict}, made in
  each of the $d - 2$ heights.

  The case $d = 3$ uses its dimension only in (2), where the point at a
  body has one height. In dimension $d$ it has $d - 2$ heights, and each
  is read at that body alone.
\end{proof}

\begin{theorem}[A cut vertex, in every dimension]
  \label{thm:pencil-d-cut}
  \uses{def:pencil-d-x0-attains, def:pencil-x0-standing}
  Let $K$ be infinite and $d \ge 3$, and let $G$ satisfy the standing
  hypotheses (\cref{def:pencil-x0-standing}), with
  $V(G) = V_1 \cup V_2$, $V_1 \cap V_2 = \{v\}$ and every edge of $G$
  inside $V_1$ or inside $V_2$. If the general configurations of $G[V_1]$
  and $G[V_2]$ attain in dimension $d$
  (\cref{def:pencil-d-x0-attains}), then the general configuration of $G$
  attains in dimension $d$.
\end{theorem}
\begin{proof}
  \uses{lem:pencil-d-lifting-restrict, lem:pencil-cut-fibre,
        lem:block-rank-cut-vertex, lem:deficiency-cut-vertex,
        lem:pencil-x0-fibre-intersection, lem:pencil-x0-main-picture-open,
        lem:pencil-d-x0-one-witness}
  \emph{The picture.} For $i = 1, 2$ the attainment at $G[V_i]$ gives a
  nonzero polynomial $P_i$ in the picture coordinates
  (\cref{def:pencil-d-x0-attains}). The graph $G$ has no loop, and each
  of its closed neighbourhoods has at least three members. So some
  nonzero polynomial $P_0$ has only main pictures of $G$ as non-roots
  (\cref{lem:pencil-x0-main-picture-open}). As $K$ is infinite,
  $P_0 P_1 P_2$ has a non-root $q$, and $q$ is a main picture of $G$. For
  $i = 1, 2$ the heights $y \in L_{G[V_i]}(q)^{d-2}$ at which $G[V_i]$ has
  its target contain a nonempty Zariski-open subset of this linear space.
  So some polynomial $R_i$ in the heights is nonzero somewhere on
  $L_{G[V_i]}(q)^{d-2}$, and $G[V_i]$ has its target at every $y$ there
  with $R_i(y) \neq 0$.

  \emph{The heights.} Read at the restriction to $V_i$, $R_i$ agrees with
  a polynomial $R_i'$ in the heights of $G$ (part~(3) of
  \cref{lem:pencil-d-lifting-restrict}). Let
  $y \in L_{G[V_i]}(q)^{d-2}$ with $R_i(y) \neq 0$. At every picture the
  restriction maps $L_G(q)$ onto $L_{G[V_i]}(q)$
  (\cref{lem:pencil-cut-fibre}). So each height $y^j$ is the restriction
  of some $\hat y^j \in L_G(q)$. The tuple
  $\hat y = (\hat y^1, \dots, \hat y^{d-2})$ lies in $L_G(q)^{d-2}$ and
  restricts to $y$, so $R_i'(\hat y) = R_i(y) \neq 0$. Thus $R_1'$ and
  $R_2'$ are each nonzero somewhere on the linear space $L_G(q)^{d-2}$,
  and some $z$ in it makes both nonzero
  (\cref{lem:pencil-x0-fibre-intersection}). The restriction $z_i$ of $z$
  to $V_i$ lies in $L_{G[V_i]}(q)^{d-2}$ (part~(1) of
  \cref{lem:pencil-d-lifting-restrict}), and $R_i(z_i) = R_i'(z) \neq 0$.
  So $G[V_i]$ has its target at $(q, z_i)$, and hence at $(q, z)$, by
  part~(2) of the same lemma.

  \emph{The count.} At $(q, z)$ the rank of $G$ is the sum of the ranks of
  $G[V_1]$ and $G[V_2]$ (\cref{lem:block-rank-cut-vertex}), a lemma on
  body-hinge frameworks in every dimension. As $G$ has no loop,
  $\operatorname{def}_d(G) = \operatorname{def}_d(G[V_1]) +
  \operatorname{def}_d(G[V_2])$ (\cref{lem:deficiency-cut-vertex}, at
  $D$). As $|V_1| + |V_2| = |V(G)| + 1$,
  \[
    \operatorname{tgt}_d(G[V_1]) + \operatorname{tgt}_d(G[V_2])
    \;=\; D(|V_1| + |V_2| - 2) - \operatorname{def}_d(G[V_1])
    - \operatorname{def}_d(G[V_2])
    \;=\; \operatorname{tgt}_d(G).
  \]
  So $(q, z)$ has rank $\operatorname{tgt}_d(G)$, with $q$ a main picture
  of $G$ and $z \in L_G(q)^{d-2}$. As $G$ is connected, it has a body, and
  \cref{lem:pencil-d-x0-one-witness} applies.

  At $d = 3$ this is \cref{thm:pencil-x0-cut}. There the dimension enters
  in two places: the height is one vector of $L_G(q)$, and the target is
  $6(|V(G)| - 1) - \operatorname{def}_3(G)$. In dimension $d$ the heights
  form a tuple, which restricts and extends height by height, since
  $L_G(q)^{d-2}$ is the product of $d - 2$ copies of $L_G(q)$. The two
  open conditions are polynomials on this one linear space, so they still
  meet. The rank identity and the deficiency identity hold at every $D$,
  so the count closes with $D$ in place of $6$.
\end{proof}

\begin{theorem}[A chain of bridges, in every dimension]
  \label{thm:pencil-d-bridge}
  \uses{def:pencil-d-x0-attains, def:pencil-x0-standing}
  Let $K$ be infinite and $d \ge 3$, and let $G$ satisfy the standing
  hypotheses. Let $V(G)$ be the disjoint union of $V_1$, $V_2$ and path
  bodies $x_1, \dots, x_k$ ($k \ge 0$), with a path
  $a - x_1 - \cdots - x_k - b$ from $a \in V_1$ to $b \in V_2$, every
  other edge of $G$ lying inside $V_1$ or inside $V_2$. If the general
  configurations of $G[V_1]$ and $G[V_2]$ attain in dimension $d$, then
  the general configuration of $G$ attains in dimension $d$.
\end{theorem}
\begin{proof}
  \uses{lem:pencil-d-lifting-restrict, lem:pencil-bridge-fibre,
        lem:block-rank-cut, lem:cut-edge-decomposition,
        lem:pencil-d-config-realization, def:panel-support-extensor,
        lem:pencil-x0-fibre-intersection, lem:pencil-x0-main-picture-open,
        lem:pencil-d-x0-one-witness}
  We find $q$ and $z$ as in the proof of \cref{thm:pencil-d-cut}, with
  \cref{lem:pencil-bridge-fibre} for the restrictions, applied to each
  height. For $G[V_2]$ we apply it from $b$, along the reversed path. So
  $q$ is a main picture of $G$, $z \in L_G(q)^{d-2}$, and $G[V_1]$ and
  $G[V_2]$ have their targets at $(q, z)$. As at $d = 3$, we count the
  ranks and the deficiencies along the whole path at once. Applying the
  step one bridge at a time would ask a graph whose only body is a path
  body to attain, and no picture of such a graph is admissible.

  Every hinge of the configuration $(q, z)$ is nonzero. Indeed, $q$ is
  admissible and $z \in L_G(q)^{d-2}$, so the points $b_u$ and $b_v$ at
  the ends of each edge are linearly independent
  (\cref{lem:pencil-d-config-realization}). Then
  $C(b_u, b_v) \neq 0$ (\cref{def:panel-support-extensor}).

  Let $x_0 = a$ and $x_{k+1} = b$, and write
  $S_j = V_1 \cup \{x_1, \dots, x_j\}$ for $0 \le j \le k$. Each $S_j$
  contains $a$ and not $b$. An edge of $G$ inside $V_1$ lies inside
  $S_j$, and an edge inside $V_2$ avoids $S_j$. The path edge
  $x_i x_{i+1}$ has exactly one end in $S_j$ when $i = j$, and otherwise
  none or two. So $x_j x_{j+1}$ is the only edge of $G$ leaving $S_j$.

  \emph{The ranks}, all at $(q, z)$. The edge $x_k x_{k+1}$ is the only
  edge between $S_k$ and $V(G) \setminus S_k = V_2$. So the rank of $G$
  is at least the sum of the ranks of $G[S_k]$ and $G[V_2]$ and $D - 1$
  (\cref{lem:block-rank-cut}). For $j < k$ the edge $x_j x_{j+1}$ joins
  $S_j$ to $x_{j+1}$, so the rank of $G[S_{j+1}]$ is at least the rank of
  $G[S_j]$ plus $D - 1$, by the same lemma. Hence the rank of $G$ is at
  least
  \[
    \operatorname{tgt}_d(G[V_1]) + \operatorname{tgt}_d(G[V_2])
    + (D - 1)(k + 1).
  \]

  \emph{The deficiencies.} The graph $G$ has no loop. For $j < k$ the only
  edge leaving $S_j$ joins it to $x_{j+1}$, so
  $\operatorname{def}_d(G[S_{j+1}]) = \operatorname{def}_d(G[S_j]) + 1$
  (part~(3) of \cref{lem:cut-edge-decomposition}). The single edge between
  $S_k$ and $V_2$ gives
  \[
    \operatorname{def}_d(G) \;=\; \operatorname{def}_d(G[S_k])
    + \operatorname{def}_d(G[V_2]) + D - (D - 1)
  \]
  (part~(1) of the same lemma). So
  $\operatorname{def}_d(G) = \operatorname{def}_d(G[V_1]) +
  \operatorname{def}_d(G[V_2]) + k + 1$. As
  $|V(G)| - 1 = (|V_1| - 1) + (|V_2| - 1) + (k + 1)$, the term
  $D(|V(G)| - 1)$ exceeds the sum of those of the two sides by $D(k + 1)$,
  and
  \[
    \operatorname{tgt}_d(G) \;=\; \operatorname{tgt}_d(G[V_1]) +
    \operatorname{tgt}_d(G[V_2]) + (D - 1)(k + 1).
  \]
  So $(q, z)$ has rank at least $\operatorname{tgt}_d(G)$. As $G$ has a
  body, \cref{lem:pencil-d-x0-one-witness} applies.

  At $d = 3$ this is \cref{thm:pencil-x0-bridge}. There the dimension
  enters in the rank that a bridge adds, $D - 1 = 5$, and in the target,
  through $6(|V(G)| - 1)$. The rank addition at a cut edge and the
  deficiency at a cut of one edge hold at every $D$, so in dimension $d$
  each bridge adds $D - 1$ to the rank bound and to the target alike. The
  hinges are nonzero in every dimension. The first, second and last
  coordinates of the points of two adjacent bodies are their homogeneous
  picture points, which are distinct with last coordinate $1$, and so
  independent. The heights restrict and extend height by height, as at a
  cut vertex.
\end{proof}

\subsection{Contraction at an additive core in every dimension}
\label{sec:pencil-d-contract}

This section proves the contraction step of the induction in dimension
$d$ (\cref{thm:pencil-d-contract-additive}), and its case at a core of
planar deficiency zero (\cref{cor:pencil-d-contract-rigid}). At $d = 3$
they are \cref{thm:pencil-x0-contract-additive} and
\cref{cor:pencil-x0-contract-rigid}, and we keep their proof
(\cref{sec:main-component-contract}). The core $H = G[W]$ is additive in
dimension $d$ (\cref{sec:pencil-d-induction}): it is rigid in dimension
$d$, and the planar deficiencies of $H$ and of the contraction $G/H$ add
up to at most that of $G$. The picture is chosen as at a cut vertex
(\cref{sec:pencil-d-cut}). The count is the block-triangular rank bound
of KT eq.~(6.3) (\cref{lem:block-rank-contract}). The heights are
found along the contraction curve $q(t)$, on which the core shrinks to
the point of a body $r \in W$, with the rescaled lifting system $M(t)$
(\cref{def:pencil-contract-lifting-system}).

The system $M(t)$ concerns one height, and in dimension $d$ a
configuration along the curve takes one vector of unknowns of $M(t)$ for
each of its $d - 2$ heights. The bounds on the kernel of $M(t)$ are
planar, and additivity enters only there. We state them first, for one
height (\cref{lem:pencil-contract-additive-kernel}): the kernel does not
jump at $t = 0$, and the core heights of the solutions of $M(0)$ are all
of $L_H(q)$. Two conditions depend on all $d - 2$ heights at once: the
attainment at $H$, and the rank of the rows of the edges outside the core
near the collapsed placement
(\cref{lem:pencil-contract-degenerate-rank}). Both are open, so they
still meet in one linear space, that of the tuples of solutions of
$M(0)$.

\begin{lemma}[The kernel at the end of the curve, at an additive core]
  \label{lem:pencil-contract-additive-kernel}
  \uses{def:pencil-contract-lifting-system, def:pencil-lifting-space,
        def:pencil-admissible-picture, def:rigid-contraction,
        def:D-deficiency}
  Let $K$ be infinite, let $r \in W \subseteq V(G)$, and let $G/H$ be the
  contraction of $H = G[W]$ to $r$ (\cref{def:rigid-contraction}), with
  $\operatorname{def}_2(H) + \operatorname{def}_2(G/H) \le
  \operatorname{def}_2(G)$. Let $q$ be a picture of $G$ that is
  admissible for $G$, for $H$ and for $G/H$, with
  $\dim L_H(q) = 3 + \operatorname{def}_2(H)$ and
  $\dim L_{G/H}(q) = 3 + \operatorname{def}_2(G/H)$. Let $M(t)$ be the
  rescaled lifting system at $q$
  (\cref{def:pencil-contract-lifting-system}), and let $\rho$ send a
  vector $(z, h)$ of its unknowns to the height that agrees with $z$ on
  $W$ and vanishes elsewhere. Then
  $\dim \ker M(0) = 3 + \operatorname{def}_2(G)$, and $\rho$ maps
  $\ker M(0)$ onto $L_H(q)$.
\end{lemma}
\begin{proof}
  \uses{lem:pencil-polynomial-section, lem:pencil-contract-lifting-kernel,
        lem:pencil-lifting-space-deficiency,
        lem:pencil-contract-kernel-bound}
  We first bound $\dim \ker M(0)$ below, at one parameter $t \neq 0$. The
  polynomial section through the zero solution at $t = 0$
  (\cref{lem:pencil-polynomial-section}, with the parameter $t$) gives a
  polynomial $E$ in $t$ with $E(0) \neq 0$ and
  $\dim \ker M(t) \le \dim \ker M(0)$ wherever $E(t) \neq 0$. As $q$ is
  admissible for $G$, some polynomial $A$ in the picture coordinates has
  $A(q) \neq 0$, and its non-roots are admissible pictures of $G$. For
  $A$ take the product, over the edges $uv$, of one of $x_u - x_v$ and
  $y_u - y_v$ that is nonzero at $q$, and, over the bodies, of one
  determinant of three homogeneous picture points of the closed
  neighbourhood that are independent at $q$ (as after
  \cref{def:pencil-admissible-picture}). Read along the curve, $A(q(t))$
  is a polynomial in $t$, nonzero at $t = 1$ as $q(1) = q$. As $K$ is
  infinite, some $t \neq 0$ has $E(t) \neq 0$ and $A(q(t)) \neq 0$. Then
  $q(t)$ is admissible for $G$, and $G$ has the body $r$, so
  \[
    3 + \operatorname{def}_2(G) \;\le\; \dim L_G(q(t)) \;\le\;
    \dim \ker M(t) \;\le\; \dim \ker M(0)
  \]
  (\cref{lem:pencil-lifting-space-deficiency},
  \cref{lem:pencil-contract-lifting-kernel}).

  Above, the kernel is bounded through its core heights. As $q$ is
  admissible for $H$ and for $G/H$, we may use both parts of
  \cref{lem:pencil-contract-kernel-bound}. With the dimensions of $L_H(q)$
  and $L_{G/H}(q)$, and additivity, they give
  \[
    \dim \ker M(0) + 3 \;\le\; \dim \rho(\ker M(0)) + \dim L_{G/H}(q)
    \;\le\; \dim L_H(q) + \dim L_{G/H}(q) \;\le\;
    6 + \operatorname{def}_2(G).
  \]
  So $\dim \ker M(0) = 3 + \operatorname{def}_2(G)$, and every inequality
  of the second chain is an equality. Hence $\rho(\ker M(0))$, which lies
  in $L_H(q)$ by part~(1) of \cref{lem:pencil-contract-kernel-bound}, has
  the dimension of $L_H(q)$, and the two are equal.

  At $d = 3$ this argument is part of the proof of
  \cref{thm:pencil-x0-contract-additive}, with the main-picture polynomial
  of $G$ in place of $A$. The lemma concerns one height and no point of a
  configuration. In dimension $d$ each of the $d - 2$ heights is the
  heights of a solution of the same system $M(t)$.
\end{proof}

For a tuple $x = (x^1, \dots, x^{d-2})$ of vectors of unknowns of
$M(t)$, the \emph{configuration of $x$ at $t$} is the configuration
$(q(t), z(t))$ in dimension $d$, where
$z(t) = (z^1(t), \dots, z^{d-2}(t))$ and $z^i(t)$ is the heights of $x^i$
at $t$ (\cref{def:pencil-contract-lifting-system}). Its points are
polynomials in $t$ and in the unknowns. We write $\ker M(t)^{d-2}$ for the
tuples whose every member solves $M(t)$, a linear subspace of the space
of tuples, and $\rho(x) = (\rho(x^1), \dots, \rho(x^{d-2}))$, with $\rho$
as in \cref{lem:pencil-contract-additive-kernel}.

\begin{theorem}[Contraction at an additive core, in every dimension]
  \label{thm:pencil-d-contract-additive}
  \uses{def:pencil-d-x0-attains, def:pencil-x0-standing, def:cut-edges-2ec,
        def:rigid-contraction, def:D-deficiency}
  Let $K$ be infinite and $d \ge 3$, and let $G$ satisfy the standing
  hypotheses (\cref{def:pencil-x0-standing}) and be $2$-edge-connected.
  Let $W \subsetneq V(G)$ with $|W| \ge 2$ and
  $\operatorname{def}_d(G[W]) = 0$, such that no body outside $W$ is
  adjacent to two bodies of $W$, and let $G/H$ be the contraction of
  $H = G[W]$ to a body $r \in W$ (\cref{def:rigid-contraction}). Suppose
  $\operatorname{def}_2(H) + \operatorname{def}_2(G/H) \le
  \operatorname{def}_2(G)$. If the general configurations of $H$ and of
  $G/H$ attain in dimension $d$ (\cref{def:pencil-d-x0-attains}), then
  the general configuration of $G$ attains in dimension $d$.
\end{theorem}
\begin{proof}
  \uses{lem:pencil-contract-standing-rigid, lem:pencil-contract-standing,
        lem:pencil-contraction-deficiency, thm:pencil-jj-equality,
        lem:pencil-x0-main-picture-open, lem:pencil-lifting-space-affine,
        lem:pencil-d-rank-scale-shift, lem:pencil-contract-additive-kernel,
        lem:pencil-d-config-realization, def:panel-support-extensor,
        lem:pencil-contract-limit, lem:pencil-contract-degenerate-rank,
        lem:pencil-x0-fibre-intersection, lem:pencil-polynomial-section,
        lem:pencil-contract-lifting-kernel,
        lem:pencil-lifting-space-deficiency,
        lem:pencil-contract-kernel-bound, lem:pencil-d-lifting-restrict,
        lem:pencil-d-contract-magnified-rank, lem:block-rank-contract,
        lem:reduction-measure, lem:pencil-d-x0-one-witness}
  We follow the proof of \cref{thm:pencil-x0-contract-additive}, with
  $d - 2$ heights.

  \emph{The pieces.} The core $H$ satisfies the standing hypotheses
  (\cref{lem:pencil-contract-standing-rigid}, at $d$). By
  \cref{lem:pencil-contract-standing}, $G/H$ is simple and satisfies them
  too. As $H$ is a proper rigid subgraph of $G$ at $D$, contracting it
  does not change the deficiency at $D$
  (\cref{lem:pencil-contraction-deficiency}), so
  $\operatorname{def}_d(G/H) = \operatorname{def}_d(G)$. The target of
  $H$ is $\operatorname{tgt}_d(H) = D(|W| - 1)$.

  \emph{The picture.} As $H$ and $G/H$ satisfy the standing hypotheses,
  Jackson and Jord\'an's equality applies to both
  (\cref{thm:pencil-jj-equality}). Take a picture $q$ off the zero sets of
  five nonzero polynomials: those of the attainment in dimension $d$ and
  of the equality at $H$ and at $G/H$, and the main-picture polynomial of
  $G$ (\cref{lem:pencil-x0-main-picture-open}). Then $q$ is a main picture
  of $G$, admissible for $H$ and for $G/H$, with
  $\dim L_H(q) = 3 + \operatorname{def}_2(H)$ and
  $\dim L_{G/H}(q) = 3 + \operatorname{def}_2(G/H)$. Some polynomial $R$
  in the heights of $H$ is nonzero somewhere on $L_H(q)^{d-2}$, and $H$
  has its target at every $y \in L_H(q)^{d-2}$ with $R(y) \neq 0$. The
  graph $G/H$ has attaining heights $y'$ over $q$. Let
  $z'^i = y'^i - y'^i_r$ for each $i$. A constant height lies in
  $\mathrm{Aff}(q) \subseteq L_{G/H}(q)$
  (\cref{lem:pencil-lifting-space-affine}), so
  $z' \in L_{G/H}(q)^{d-2}$. The configuration $(q, z')$ has the rank of
  $(q, y')$ (\cref{lem:pencil-d-rank-scale-shift}, with $T$ the
  identity). So $z'$ is an attaining height of $G/H$, with $z'^i_r = 0$
  for every $i$.

  \emph{The kernel at the end of the curve.} Let $M(t)$ be the rescaled
  lifting system of $G$ at $q$, for $W$ and $r$. By
  \cref{lem:pencil-contract-additive-kernel},
  $\dim \ker M(0) = 3 + \operatorname{def}_2(G)$, and $\rho$ maps
  $\ker M(0)$ onto $L_H(q)$. Applied to each member, $\rho$ therefore maps
  $\ker M(0)^{d-2}$ onto $L_H(q)^{d-2}$.

  \emph{Two open conditions.} First, $R(\rho(x))$ is a polynomial in the
  tuple $x$. It is nonzero somewhere on $\ker M(0)^{d-2}$: take
  $y \in L_H(q)^{d-2}$ with $R(y) \neq 0$, and $x \in \ker M(0)^{d-2}$
  with $\rho(x) = y$.

  Second, let $b'$ be the points of $(q, z')$. As $q$ is admissible for
  $G/H$ and $z' \in L_{G/H}(q)^{d-2}$, the points $b'_u$ and $b'_v$ at the
  ends of each edge of $G/H$ are linearly independent
  (\cref{lem:pencil-d-config-realization}). So every hinge
  $C(b'_u, b'_v)$ of $G/H$ is nonzero (\cref{def:panel-support-extensor}),
  and the rank of $G/H$ at the normals $b'$ is
  $\operatorname{tgt}_d(G/H)$. Let $Q$ be the polynomial of
  \cref{lem:pencil-contract-degenerate-rank} for the normals $b'$ and
  $N = \operatorname{tgt}_d(G/H)$. Read at the configuration of $x$ at
  $t = 0$, it is a polynomial in $x$. For each $i$ there is a solution
  $x'^i$ of $M(0)$ whose heights agree with $z'^i$ off $W$
  (\cref{lem:pencil-contract-limit}). At $t = 0$ every body of $W$ has the
  picture point $q_r$ and all heights zero. So the configuration of
  $x' = (x'^1, \dots, x'^{d-2})$ at $t = 0$ puts every body of $W$ at
  $b'_r$, as every $z'^i_r$ vanishes, and every other body $v$ at $b'_v$.
  This is the collapsed placement of $b'$, so $Q$ is nonzero at $x'$. Some
  $x_0 \in \ker M(0)^{d-2}$ makes both polynomials nonzero
  (\cref{lem:pencil-x0-fibre-intersection}).

  \emph{The curve.} For each $i$, apply
  \cref{lem:pencil-polynomial-section} to $M(t)$ at $t = 0$ and $x_0^i$.
  It gives a polynomial $D_i$ in $t$ with $D_i(0) \neq 0$, and a vector
  $X^i(t)$ of polynomials in $t$ with $X^i(0) = x_0^i$. For all but
  finitely many $t$, $q(t)$ is a main picture of $G$, as the main-picture
  polynomial of $G$ along the curve is nonzero at $t = 1$, where
  $q(1) = q$. At every $t \neq 0$ with $q(t)$ a main picture of $G$ the
  kernel does not jump:
  \[
    \dim \ker M(0) \;=\; 3 + \operatorname{def}_2(G) \;\le\;
    \dim L_G(q(t)) \;\le\; \dim \ker M(t)
  \]
  (\cref{lem:pencil-lifting-space-deficiency},
  \cref{lem:pencil-contract-lifting-kernel}). So at every such $t$ with
  $D_i(t) \neq 0$ for every $i$, the tuple
  $X(t) = (X^1(t), \dots, X^{d-2}(t))$ lies in $\ker M(t)^{d-2}$, and the
  heights $z(t)$ of its configuration at $t$ lie in $L_G(q(t))^{d-2}$
  (\cref{lem:pencil-contract-lifting-kernel}, at each height).

  \emph{A good parameter.} Take $t$ off the zero sets of $t$, of
  $D_1, \dots, D_{d-2}$, of the main-picture polynomial of $G$ along the
  curve, and of the two polynomials of the open conditions along $X(t)$:
  $R(\rho(X(t)))$, and $Q$ at the configuration of $X(t)$ at $t$. These
  two are nonzero at $t = 0$, where they are the open conditions at
  $x_0$. Each $\rho(X^i(t))$ lies in $L_H(q)$ (part~(1) of
  \cref{lem:pencil-contract-kernel-bound}), so $\rho(X(t))$ is an
  attaining height of $H$ over $q$, and $H$ has rank $D(|W| - 1)$ at
  $(q, \rho(X(t)))$. Write $X^i(t) = (\hat z^i, \hat h^i)$ and
  $\hat z = (\hat z^1, \dots, \hat z^{d-2})$. The tuple $\rho(X(t))$
  agrees with $\hat z$ on $W$ and vanishes elsewhere, so $H$ has the same
  rank at $(q, \hat z)$ (part~(2) of \cref{lem:pencil-d-lifting-restrict}).
  It has this rank at $(q(t), z(t))$ as well
  (\cref{lem:pencil-d-contract-magnified-rank}). As $Q$ is nonzero at the
  points of $(q(t), z(t))$, the rows there of the edges of $G$ not in
  $H$, with the columns of $W$ deleted, have rank at least
  $\operatorname{tgt}_d(G/H)$. By \cref{lem:block-rank-contract} the rank
  of $G$ at $(q(t), z(t))$ is at least
  \[
    D(|W| - 1) + \operatorname{tgt}_d(G/H) \;=\;
    D(|W| - 1) + D(|V(G)| - |W|) - \operatorname{def}_d(G) \;=\;
    \operatorname{tgt}_d(G),
  \]
  since $|V(G/H)| = |V(G)| - |W| + 1$ (\cref{lem:reduction-measure}) and
  $\operatorname{def}_d(G/H) = \operatorname{def}_d(G)$. So $q(t)$ is a
  main picture of $G$, $z(t) \in L_G(q(t))^{d-2}$, and $(q(t), z(t))$ has
  rank at least $\operatorname{tgt}_d(G)$. As $G$ has a body,
  \cref{lem:pencil-d-x0-one-witness} applies.

  At $d = 3$ this is \cref{thm:pencil-x0-contract-additive}. There the
  dimension enters in four places. First, the core is rigid at $D = 6$.
  In dimension $d$ it is rigid at $D$, and the three facts drawn from its
  rigidity hold at every $D$: the core satisfies the standing hypotheses,
  contracting it keeps the deficiency, and its target is $D(|W| - 1)$.
  Additivity stays planar. It is used only in the kernel at the end of
  the curve, which is planar (\cref{lem:pencil-contract-additive-kernel}).
  Second, the height is one vector. In dimension $d$ each of the $d - 2$
  heights is the heights of its own vector of unknowns of $M(t)$. The
  system, its solutions at the end of the curve
  (\cref{lem:pencil-contract-limit}) and the polynomial section act on
  each height separately, and the shift that makes $z'_r = 0$ is made in
  each height. Third, the attainment at $H$ and the rank of the rows
  outside the core each depend on all $d - 2$ heights at once. Both are
  open conditions, so they still meet, in the linear space
  $\ker M(0)^{d-2}$. Fourth, the rank of the core along the curve is
  compared with its rank over $q$ through an invertible linear map of
  $K^{d+1}$ in place of one of $K^4$, and the count closes with $D$ in
  place of $6$. The block-triangular rank bound and the rank of the
  surviving rows near the collapsed placement
  (\cref{lem:block-rank-contract},
  \cref{lem:pencil-contract-degenerate-rank}) hold for panel-hinge
  frameworks in every dimension, with the same proofs.
\end{proof}

\begin{corollary}[Contraction at a core of planar deficiency zero, in every dimension]
  \label{cor:pencil-d-contract-rigid}
  \uses{def:pencil-d-x0-attains, def:pencil-x0-standing, def:cut-edges-2ec,
        def:rigid-contraction, def:D-deficiency}
  Let $K$ be infinite and $d \ge 3$, and let $G$ satisfy the standing
  hypotheses and be $2$-edge-connected. Let $W \subsetneq V(G)$ with
  $|W| \ge 2$ and $\operatorname{def}_2(G[W]) = 0$, such that no body
  outside $W$ is adjacent to two bodies of $W$, and let $G/H$ be the
  contraction of $H = G[W]$ to a body of $W$. If the general
  configuration of $G/H$ attains in dimension $d$, then the general
  configuration of $G$ attains in dimension $d$.
\end{corollary}
\begin{proof}
  \uses{thm:pencil-d-contract-additive, lem:pencil-contract-standing-rigid,
        lem:deficiency-rigid-higher, cor:pencil-d-jj-flat,
        lem:pencil-contraction-deficiency}
  The core $H$ satisfies the standing hypotheses
  (\cref{lem:pencil-contract-standing-rigid}, at $d = 2$). As $H$ has a
  body and $\operatorname{def}_2(H) = 0$, also
  $\operatorname{def}_d(H) = 0$ (\cref{lem:deficiency-rigid-higher}). So
  $W$ is rigid in dimension $d$, and the two deficiencies of $H$ agree.
  Hence the general configuration of $H$ attains in dimension $d$
  (\cref{cor:pencil-d-jj-flat}). Contracting a core of planar deficiency
  zero does not change the planar deficiency
  (\cref{lem:pencil-contraction-deficiency}, at $d = 2$). So
  $\operatorname{def}_2(H) + \operatorname{def}_2(G/H) =
  \operatorname{def}_2(G)$, and \cref{thm:pencil-d-contract-additive}
  applies.

  At $d = 3$ this is \cref{cor:pencil-x0-contract-rigid}. There the
  dimension enters in two places: the bound
  $\operatorname{def}_3(H) \le \operatorname{def}_2(H)$ makes the core
  rigid, and the case of equal deficiencies at $d = 3$ makes its general
  configuration attain. In dimension $d$ these are
  \cref{lem:deficiency-rigid-higher} and \cref{cor:pencil-d-jj-flat}. The
  standing hypotheses of the core and the planar deficiency of the
  contraction do not depend on $d$.
\end{proof}

\subsection{Splitting off a body of degree two in every dimension}
\label{sec:pencil-d-splitoff}

This section proves the split-off step of the induction in dimension $d$
(\cref{thm:pencil-d-splitoff}). At $d = 3$ it is
\cref{thm:pencil-x0-splitoff}, and we keep its proof
(\cref{sec:main-component-splitoff}). Let $x$ be a body of degree two
whose neighbours $a$ and $b$ are not adjacent, let
$V_1 = V(G) \setminus \{x\}$, and let $G''$ be $G[V_1]$ with a new edge
$ab$. The step assumes
\[
  \delta_d \;=\; \operatorname{def}_d(G[V_1]) -
  \operatorname{def}_d(G[V_1]; a, b) \;\ge\; D - 1,
\]
where the step at $d = 3$ assumes $\delta_3 \ge 5$. We put $x$ back on
the line through the points of $a$ and $b$, where it raises the rank by
exactly $D - 1$ (\cref{lem:pencil-splitoff-special-rank}). Then we move
its picture along a line to a general point, and the heights along a
line of solutions of the lifting system of $G[V_1]$.

The lemmas of \cref{sec:main-component-splitoff} serve in dimension $d$
unchanged. Three of them are planar: the alternative that places the
special position (\cref{lem:pencil-splitoff-flexes}), the line of
solutions (\cref{lem:pencil-splitoff-curve}) and the height at the ear
body (\cref{lem:pencil-splitoff-ear-extension}). In dimension $d$ we
apply them to each of the $d - 2$ heights. The other two concern
panel-hinge frameworks with normals in any dimension: the rank at the
special position (\cref{lem:pencil-splitoff-special-rank}) and the rank
along a polynomial curve (\cref{lem:pencil-curve-limit}).

At a chain $a - x - b$ on $V_1$ with one interior body, the last clause
of \cref{def:pencil-d-usable-chain} is the bound $\delta_d \ge D - 1$.
This bound holds when no set of at least two bodies containing $x$ is
rigid in dimension $d$ (\cref{lem:deficiency-one-body-chain-d}). In a
sparse graph the ends $a$ and $b$ of such a chain are not adjacent, by
part~(1) of \cref{lem:deficiency-one-body-chain}, which is planar.

\begin{theorem}[Splitting off a body of degree two, in every dimension]
  \label{thm:pencil-d-splitoff}
  \uses{def:pencil-d-x0-attains, def:pencil-x0-standing, def:open-ear,
        def:graph-operations, def:D-deficiency, def:deficiency-merged}
  Let $K$ be infinite and $d \ge 3$, with $D = \binom{d+1}{2}$, and let
  $G$ satisfy the standing hypotheses (\cref{def:pencil-x0-standing}).
  Let $a - x - b$ be an open ear of $G$ on $V_1$ with one interior body
  (\cref{def:open-ear}), whose ends $a$ and $b$ are not adjacent. Let
  $G''$ be the splitting-off $G_x^{ab}$ (\cref{def:graph-operations})
  whose new edge $ab$ takes the label of the edge $ax$. Suppose
  $\operatorname{def}_d(G[V_1]; a, b) + D - 1 \le
  \operatorname{def}_d(G[V_1])$ (\cref{def:deficiency-merged}). If the
  general configuration of $G''$ attains in dimension $d$
  (\cref{def:pencil-d-x0-attains}), then the general configuration of
  $G$ attains in dimension $d$.
\end{theorem}
\begin{proof}
  \uses{lem:splitoff-deficiency-merged, lem:deficiency-ear,
        lem:splitoff-deficiency-reuse, thm:pencil-jj-equality,
        lem:pencil-lifting-space-deficiency,
        lem:pencil-x0-main-picture-open, def:pencil-lifting-system,
        lem:pencil-splitoff-flexes, lem:pencil-splitoff-curve,
        lem:pencil-weighted-lifting-system-local,
        lem:pencil-splitoff-special-rank, lem:pencil-rank-congr,
        lem:pencil-d-config-realization, def:panel-support-extensor,
        lem:pencil-curve-limit, lem:pencil-splitoff-ear-extension,
        lem:pencil-d-x0-one-witness}
  We follow the proof of \cref{thm:pencil-x0-splitoff}, with $d - 2$
  heights.

  \emph{The counts.} The new edge $ab$ lowers by $D - 1$ the value of
  every partition of $V_1$ that separates $a$ and $b$. By hypothesis the
  other partitions have value at most
  $\operatorname{def}_d(G[V_1]) - (D - 1)$. So
  $\operatorname{def}_d(G'') + D - 1 \le \operatorname{def}_d(G[V_1])$
  (\cref{lem:splitoff-deficiency-merged}, at $D$). The ear $a - x - b$
  has one interior body, so
  $\operatorname{def}_d(G[V_1]) \le \operatorname{def}_d(G) + D - 2$
  (\cref{lem:deficiency-ear}, at $k = 1$). Together these give
  $\operatorname{def}_d(G'') + 1 \le \operatorname{def}_d(G)$. Also
  $\operatorname{def}_2(G'') \le \operatorname{def}_2(G)$
  (\cref{lem:splitoff-deficiency-reuse}, at $D = 3$).

  \emph{The picture.} As $a$ and $b$ are not adjacent, $G''$ is simple,
  and each of its closed neighbourhoods has as many members as in $G$,
  at least three. So Jackson and Jord\'an's equality applies to $G''$
  (\cref{thm:pencil-jj-equality}). Take a picture $q$ off the zero sets
  of three nonzero polynomials: that of the attainment at $G''$ in
  dimension $d$, that of the equality at $G''$, and the main-picture
  polynomial of $G$ (\cref{lem:pencil-x0-main-picture-open}). Then $q$
  is admissible for $G''$ and a main picture of $G$, and some
  $z'' \in L_{G''}(q)^{d-2}$ gives $G''$ its target. Also
  \[
    \dim L_{G''}(q) \;=\; 3 + \operatorname{def}_2(G'') \;\le\;
    3 + \operatorname{def}_2(G) \;\le\; \dim L_G(q)
  \]
  (\cref{lem:pencil-lifting-space-deficiency}). For each $i$ let
  $y_0^i = (z''^i, h^i)$ be a solution of the lifting system of $G''$ at
  $q$ with heights $z''^i$ (\cref{def:pencil-lifting-system}). Each
  closed neighbourhood in $G[V_1]$ lies inside the one in $G''$, so
  $y_0^i$ is a solution for $G[V_1]$. As $ab$ is an edge of $G''$, the
  affine function $h^i_a - h^i_b$ vanishes at $(q_a, 1)$ and at
  $(q_b, 1)$.

  \emph{The special position.} If some solution for $G[V_1]$ has
  $h_a - h_b$ nonzero at $(q_a, 1)$ or at $(q_b, 1)$, then for all but
  at most one $s$ it is nonzero at $(1 - s)(q_a, 1) + s\,(q_b, 1)$.
  Otherwise every solution for $G[V_1]$ has $h_a = h_b$
  (\cref{lem:pencil-splitoff-flexes}). This alternative reads the
  solutions for $G[V_1]$ at $q$, and not the heights $z''$. As $K$ is
  infinite, we may therefore fix one $s \neq 0, 1$ accordingly for all
  $d - 2$ heights. Put $q^0_x = (1 - s)\,q_a + s\,q_b$, and let
  $\varphi_0(z, h) = (h_a - h_b) \cdot (q^0_x, 1)$ and
  $\psi(z, h) = (h_a - h_b) \cdot (q_x - q^0_x, 0)$. Then $\varphi_0$ is
  nonzero somewhere on the solutions for $G[V_1]$, or $\varphi_0$ and
  $\psi$ both vanish on them. Each $\varphi_0(y_0^i)$ is zero, as
  $(q^0_x, 1)$ is an affine combination of $(q_a, 1)$ and $(q_b, 1)$.
  By \cref{lem:pencil-splitoff-curve}, applied once for each height,
  there are solutions $w^i$ for $G[V_1]$ such that
  $y^i(t) = y_0^i + t\,w^i$ satisfies
  $\varphi_0(y^i(t)) + t\,\psi(y^i(t)) = 0$ for every $t$.

  \emph{The curve.} Let $q(t)$ be $q$ with the picture of $x$ moved to
  $q^0_x + t\,(q_x - q^0_x)$, so that $q(1) = q$. As $q(t)$ agrees with
  $q$ on $V_1$, each $y^i(t)$ is a solution for $G[V_1]$ at $q(t)$
  (\cref{lem:pencil-weighted-lifting-system-local}). Write
  $y^i(t) = (u^i(t), h^i(t))$. The identity above says that
  $h^i_a(t) - h^i_b(t)$ vanishes at $(q(t)_x, 1)$. Let $z^i(t)$ agree
  with $u^i(t)$ on $V_1$ and take the value $h^i_a(t) \cdot (q(t)_x, 1)$
  at $x$, and let $z(t) = (z^1(t), \dots, z^{d-2}(t))$. The points of the
  configuration $(q(t), z(t))$ in dimension $d$ are polynomials in $t$.

  At $t = 0$ they are at the special position. On $V_1$ they are the
  points $p_v$ of $(q, z'')$, as $y^i(0) = y_0^i$. At $x$ the picture
  point is $q^0_x$, and the $i$-th height is
  \[
    h^i_a \cdot \bigl((1 - s)\,(q_a, 1) + s\,(q_b, 1)\bigr)
    \;=\; (1 - s)\,z''^i_a + s\,z''^i_b,
  \]
  as $h^i_a$ takes the values $z''^i_a$ and $z''^i_b$ at $q_a$ and
  $q_b$. The last coordinate of each point is $1 = (1 - s) + s$. So the
  point at $x$ is $p_x = (1 - s)\,p_a + s\,p_b$ in $K^{d+1}$. The points
  $p_a$ and $p_b$ are linearly independent: their last coordinates are
  $1$, and $q_a \neq q_b$, as $q$ is admissible for $G''$, in which $a$
  and $b$ are adjacent. So at $t = 0$ the rank of $G$ is that of $G''$
  at $(q, z'')$ plus $D - 1$ (\cref{lem:pencil-splitoff-special-rank},
  with normals in $K^{d+1}$, and \cref{lem:pencil-rank-congr}). This is
  $\operatorname{tgt}_d(G'') + D - 1$.

  \emph{A good parameter.} Every hinge of $G$ is nonzero at $t = 0$. At
  the edges $ax$ and $xb$ the picture point $q^0_x$ differs from $q_a$
  and from $q_b$, as $s \neq 0, 1$. So the points at the two ends, whose
  last coordinates are $1$, are linearly independent. Every other edge of
  $G$ is an edge of $G''$, and the points of $(q, z'')$ at its ends are
  linearly independent (\cref{lem:pencil-d-config-realization}). A hinge
  $C(p_u, p_v)$ with $p_u$ and $p_v$ independent is nonzero
  (\cref{def:panel-support-extensor}). By \cref{lem:pencil-curve-limit},
  the rank of $G$ at $(q(t), z(t))$ is therefore at least its value at
  $t = 0$ for all but finitely many $t$. The main-picture polynomial of
  $G$ at $q(t)$ is a polynomial in $t$, nonzero at $t = 1$, so $q(t)$ is
  a main picture of $G$ for all but finitely many $t$. Take $t$ outside
  both finite sets. Then $q(t)$ is admissible for $G$, and each $z^i(t)$
  lies in $L_G(q(t))$ (\cref{lem:pencil-splitoff-ear-extension}, at
  $q(t)$), so $z(t) \in L_G(q(t))^{d-2}$. As $|V(G)| = |V_1| + 1$ and
  $\operatorname{def}_d(G'') + 1 \le \operatorname{def}_d(G)$, the rank
  of $G$ at $(q(t), z(t))$ is at least
  \[
    D(|V_1| - 1) - \operatorname{def}_d(G'') + D - 1 \;=\;
    D(|V(G)| - 1) - \operatorname{def}_d(G'') - 1 \;\ge\;
    \operatorname{tgt}_d(G).
  \]
  As $G$ has a body, \cref{lem:pencil-d-x0-one-witness} applies, and the
  general configuration of $G$ attains in dimension $d$.

  At $d = 3$ this is \cref{thm:pencil-x0-splitoff}. There the dimension
  enters in three places. First, the count reads $\delta_3 \ge 5$, the
  five that $x$ adds to the rank at the special position, the bound
  $\operatorname{def}_3(G[V_1]) \le \operatorname{def}_3(G) + 4$ along
  the ear, and the target, through $6(|V(G)| - 1)$. In dimension $d$
  these are $D - 1$, $D - 1$, $D - 2$ and $D(|V(G)| - 1)$.
  \Cref{lem:splitoff-deficiency-merged} and \cref{lem:deficiency-ear}
  hold at every $D$, and the count closes as at $d = 3$. Second, the height is one vector. In dimension $d$ each of
  the $d - 2$ heights has its own solution $y_0^i$, its own line
  $y^i(t)$ and its own value at $x$. The alternative that fixes $s$, and
  the functionals $\varphi_0$ and $\psi$, read only the lifting system of
  $G[V_1]$ at $q$, which is planar. So one special position serves every
  height. Third, at $d = 3$ the special position puts the point of $x$ on
  the line through those of $a$ and $b$ in $K^4$. In dimension $d$ it
  lies on that line in $K^{d+1}$, since each of its heights is the value
  of an affine function at an affine combination of the picture points
  of $a$ and $b$. The rank at the special position and the rank along a
  polynomial curve are lemmas on panel-hinge frameworks with normals in
  any dimension.
\end{proof}

\subsection{Chains with few interior bodies in every dimension}
\label{sec:pencil-d-chains}

At $d = 3$, once the first five cases of the coverage fail, a usable
chain gives a step, and a core is found when no chain is usable
(\cref{sec:main-component-coverage}). This core exists in every
dimension: if no chain is usable in dimension $d$, then $G$ has a core
in dimension $d$ (\cref{thm:pencil-d-theorem-s}). But a usable chain
with one interior body, or with two at non-adjacent ends, may meet only
the conditions of the ear steps of \cref{sec:main-component-orbit}
(\cref{thm:pencil-x0-open-ear-one} and
\cref{thm:pencil-x0-open-ear-two-orbit}). Both steps place the point of
$x_1$ in $\Pi_a \cap \Pi_b$. Two three-dimensional subspaces of $K^4$
meet in a subspace of dimension at least two, while for $d \ge 5$ two
three-dimensional subspaces of $K^{d+1}$ can meet only in zero. We do
not extend these two steps, and count edges along the chains instead.

Let every chain have at most $\ell$ interior bodies, where
$3(\ell + 1) \le D$. Then $s_D(V_1) < 0$ for every chain on $V_1$
(\cref{lem:pencil-d-chain-count}). If moreover
none of the first five cases settles $G$, then $G$ has a core in
dimension $d$ unless every chain has one interior body and
$\operatorname{def}_3(G) = 0$ (\cref{thm:pencil-d-chain-or-core}). We
may take $\ell = 2$ at $d = 4$, and $\ell = 3$ at every $d \ge 5$. So
the coverage needs no step at a chain with at most two interior bodies
at $d = 4$, nor at one with at most three from $d = 5$ on.

In the coverage in dimension $d$, a chain with more than $\ell$
interior bodies is settled by an ear step with at least three interior
bodies; at $d = 3$ these are \cref{thm:pencil-x0-open-ear-three},
\cref{thm:pencil-x0-open-ear-four} and \cref{thm:pencil-x0-open-ear}. A
core is settled by the contraction (\cref{thm:pencil-d-contract-additive}).
The general configuration of $G$ attains in dimension three
(\cref{thm:pencil-x0-coverage}, \cref{lem:pencil-d-x0-attains-three}).
So when $\operatorname{def}_3(G) = \operatorname{def}_d(G)$, it attains
in dimension $d$ (\cref{lem:pencil-d-x0-attains-of-le}).

% Source: SHORT-CORE workbook (notes/attacks/short-core/workbook.md), S53 (a).
\begin{lemma}[Counting edges when every chain has few interior bodies]
  \label{lem:pencil-d-chain-count}
  \uses{def:pencil-x0-chain, def:pencil-x0-standing}
  Let $d \ge 3$ and $D = \binom{d+1}{2}$, and let $\ell \ge 1$ with
  $3(\ell + 1) \le D$. Let $G$ satisfy the standing hypotheses
  (\cref{def:pencil-x0-standing}), with $G - v$ connected for every body
  $v$ and some body of degree other than two, and let every chain of $G$
  (\cref{def:pencil-x0-chain}) have at most $\ell$ interior bodies. Then
  every chain of $G$ on $V_1$ has
  \[
    s_D(V_1) \;\le\; -(\ell + 1),
  \]
  where $s_D(V_1) = D(|V_1| - 1) - (D - 1)\,e(V_1)$ is the singleton
  value of $V_1$ (\cref{sec:main-component-sparse}).
\end{lemma}
\begin{proof}
  \uses{lem:pencil-x0-chain-exists, def:open-ear}
  Call a body of degree at least three a \emph{branch body}, and let $B$
  be the set of branch bodies. It is nonempty, as $G$ has minimum degree
  two and a body of degree other than two. Call a chain, or an edge
  joining two branch bodies, a \emph{segment}. Its \emph{length} is its
  number of edges, one more than its number of interior bodies.

  \emph{The segments.} Every body $x$ of degree two is an interior body
  of a chain (\cref{lem:pencil-x0-chain-exists}). On a chain the two
  bodies next to an interior body are distinct and adjacent to it, so
  they are its two neighbours. So the chain through $x$ is found by
  walking from $x$ in both directions through bodies of degree two up to
  the first branch body, and it is unique up to reversal. An edge with an
  end of degree two lies on the chain of that end. It lies on no other
  segment, since a segment through it has that end as an interior body.
  An edge joining two branch bodies lies on no chain, as every
  edge of a chain has an interior end. So every edge of $G$ lies on
  exactly one segment, and every body of degree two is an interior body
  of exactly one. Let the segments be $\sigma_1, \dots, \sigma_e$, of
  lengths $L_1, \dots, L_e$. Then
  \[
    |E(G)| \;=\; \sum_j L_j, \qquad
    |V(G)| \;=\; |B| + \sum_j (L_j - 1).
  \]

  \emph{The degrees.} The ends of a segment are two distinct branch
  bodies, as the ends of an open ear are distinct (\cref{def:open-ear})
  and $G$ is simple. A branch body $v$ is an interior body of no
  segment. So each edge at $v$ lies on a segment that ends at $v$. A
  segment that ends at $v$ has exactly one edge at $v$, as its bodies
  are distinct and its other end is not $v$. Hence the
  degree of $v$ is the number of segments that end at $v$, and the
  degrees of the branch bodies sum to $2e$. As each is at least three,
  $3|B| \le 2e$, that is, $e - |B| \ge e/3$.

  \emph{The count.} Let $\sigma_i$ be a chain
  $a - x_1 - \cdots - x_k - b$ on $V_1$, so that $k = L_i - 1$. By
  \cref{def:open-ear}, $V_1$ is $V(G)$ without $x_1, \dots, x_k$, and
  every edge of $G$ off the chain lies inside $V_1$. Every edge of the
  chain has an interior end. So
  \[
    |V_1| - 1 \;=\; |B| - e + \sum_{j \neq i} L_j, \qquad
    e(V_1) \;=\; \sum_{j \neq i} L_j,
  \]
  and therefore
  \[
    s_D(V_1) \;=\; \sum_{j \neq i} L_j \;-\; D\,(e - |B|).
  \]
  Every segment has at most $\ell + 1$ edges. With $e - |B| \ge e/3$
  this gives
  \[
    s_D(V_1) \;\le\; (\ell + 1)(e - 1) - \frac{D e}{3}
    \;=\; e\Bigl(\ell + 1 - \frac{D}{3}\Bigr) - (\ell + 1)
    \;\le\; -(\ell + 1),
  \]
  as $3(\ell + 1) \le D$.

  At $d = 3$ only $\ell = 1$ is allowed. A chain then has one interior
  body $x$, whose neighbours have degree at least three, and the bound
  $s_6(V(G) \setminus \{x\}) \le -2$ is the one that the count in the
  proof of part~(2) of \cref{lem:pencil-x0-sparse-count} gives. The
  dimension enters the count only through the coefficients $D$ and
  $D - 1$ of the singleton value. The bound needs $3(\ell + 1) \le D$,
  so it reaches the chains with two interior bodies at $d = 4$, where
  $D = 10$, and those with three at every $d \ge 5$, where $D \ge 15$.
\end{proof}

% Source: SHORT-CORE workbook (notes/attacks/short-core/workbook.md), S53 (b),
% with S52 and the ONE-CLASS lemma, both formalized in sec:pencil-d-deficiency.
\begin{theorem}[A chain with many interior bodies, a core, or rigidity in dimension three]
  \label{thm:pencil-d-chain-or-core}
  \uses{def:pencil-x0-chain, def:pencil-x0-standing,
        def:rigid-contraction, def:D-deficiency}
  Let $d \ge 3$ and $D = \binom{d+1}{2}$, and let $\ell \ge 1$ with
  $3(\ell + 1) \le D$. Let $G$ satisfy the standing hypotheses
  (\cref{def:pencil-x0-standing}), with $G - v$ connected for every body
  $v$, some body of degree other than two, and no planar-rigid set. Then
  $G$ has a chain (\cref{def:pencil-x0-chain}) with more than $\ell$
  interior bodies, or $G$ has a core in dimension $d$
  (\cref{sec:pencil-d-deficiency}), or every chain of $G$ has one
  interior body and
  $\operatorname{def}_3(G) = \operatorname{def}_d(G) = 0$.
\end{theorem}
\begin{proof}
  \uses{lem:deficiency-sparse, lem:pencil-d-chain-count,
        lem:pencil-d-chain-core, lem:pencil-x0-chain-exists,
        lem:pencil-d-core-positive-deficiency,
        lem:deficiency-rigid-higher}
  Suppose that every chain of $G$ has at most $\ell$ interior bodies.
  As $G$ has no planar-rigid set, it is sparse
  (\cref{lem:deficiency-sparse}).

  Let some chain $a - x_1 - \cdots - x_k - b$ on $V_1$ have $k \ge 2$.
  Then $s_D(V_1) \le -(\ell + 1) < 0$
  (\cref{lem:pencil-d-chain-count}). A chain is an open ear, so $G$ has
  a core in dimension $d$ that avoids $x_1$
  (\cref{lem:pencil-d-chain-core}).

  Otherwise every chain has one interior body, and then no two bodies
  of degree two are adjacent. Indeed, let $x$ and $y$ be adjacent bodies
  of degree two. Then $x$ is an interior body of a chain
  (\cref{lem:pencil-x0-chain-exists}). The two bodies next to $x$ on it
  are the two neighbours of $x$, so $y$ is one of them. As the ends of a
  chain have degree at least three, $y$ is an interior body too, and the
  chain has two. So if $\operatorname{def}_3(G) > 0$, then $G$ has a
  core in dimension $d$ (\cref{lem:pencil-d-core-positive-deficiency}).
  Otherwise $\operatorname{def}_3(G) = 0$, and then
  $\operatorname{def}_d(G) = 0$, as $d \ge 3$
  (\cref{lem:deficiency-rigid-higher}).

  At $d = 3$ these graphs are settled by the usable chains and
  \cref{thm:pencil-x0-theorem-s}, and the dimension enters through
  usability. A chain with three or more interior bodies is usable. A
  chain with one or two is usable under a condition on its merged
  deficiencies, at $D = 6$ or in the plane. In dimension $d$ the count
  replaces these conditions at the chains with two to $\ell$ interior
  bodies, and \cref{lem:pencil-d-core-positive-deficiency} replaces them
  when every chain has one. At $d = 3$ only $\ell = 1$ is allowed, and
  the theorem then leaves every chain with two or more interior bodies
  to the ear steps. The two lemmas that give the core hold in every
  dimension, with rigidity read at $D$ and additivity in the plane.
\end{proof}
